% Môn: Vật lý
% Lớp: 11
% Chương: Chương II. Sóng
% Bài: L11C2 Bài 8. Mô tả sóng
% Số câu: 170
% App đọc file này trực tiếp — sửa rồi Commit trên GitHub, bấm Đồng bộ GitHub .tex

% L11C2 Bài 8. Mô tả sóng

% ===== Câu 1 =====
% ID: L11C2B8-01-DS
% Mức: TH
\dangbt{Bài toán tổng hợp về các đại lượng đặc trưng của sóng}
\begin{ex}
% ID: L11C2B8-01-DS
% Mức: TH
 Một sóng có $f=10$ Hz và $v=30$ m/s.
\choiceTF
{\True Chu kì bằng $0,1\,\text{s}$.}
{\True Bước sóng bằng $3\,\text{m}$.}
{\True Trong $2\,\text{s}$ sóng truyền được $60\,\text{m}$.}
{Tăng tần số trong cùng môi trường làm tốc độ truyền sóng tăng tương ứng.}
\loigiai{
\begin{itemchoice}
\item (Đúng) Chu kì bằng $0,1\,\text{s}$. (Vì): Chu kì bằng $0,1\,\text{s}$. b) Đ: Bước sóng bằng $3\,\text{m}$. c) Đ: Trong $2\,\text{s}$ sóng truyền được $60\,\text{m}$. d) S: Tăng tần số trong cùng môi trường làm tốc độ truyền sóng tăng tương ứng
\item (Đúng) Bước sóng bằng $3\,\text{m}$.
\item (Đúng) Trong $2\,\text{s}$ sóng truyền được $60\,\text{m}$.
\item (Sai) Tăng tần số trong cùng môi trường làm tốc độ truyền sóng tăng tương ứng.
\end{itemchoice}
}
\end{ex}



% ID: L11C2B8-01-TLN
% Mức: TH
\begin{ex}
% ID: L11C2B8-01-TLN
% Mức: TH
 Sóng có tốc độ $12\,\text{m}$ /s và tần số $4\,\text{Hz}$. Bước sóng tính theo m bằng bao nhiêu?
\shortans{3}
\loigiai{
Xác định đúng dữ kiện và đơn vị, áp dụng hệ thức của dạng bài rồi tính được kết quả 3.
}
\end{ex}

% ===== Câu 4 =====

\begin{ex}
% ID: L11C2B8-01-TN
% Mức: NB
Một sóng có tần số $5,\text{Hz}$ và tốc độ $20,\text{m/s}$. Bước sóng bằng
\choice
{\True $4,\text{m}$}
{$100,\text{m}$}
{$0,25,\text{m}$}
{$15,\text{m}$}
\loigiai{
Ta có hệ thức giữa tốc độ truyền sóng, tần số và bước sóng:

$$
v = f\lambda.
$$

Suy ra

$$
\lambda = \frac{v}{f} = \frac{20}{5} = 4\,\text{m}.
$$

Vậy đáp án đúng là $\boxed{4,\text{m}}$.
}
\end{ex}


% ===== Câu 5 =====
% ID: L11C2B8-02-DS
% Mức: TH
\begin{ex}
% ID: L11C2B8-02-DS
% Mức: TH
Trên ảnh sóng, 6 đỉnh liên tiếp trải trên đoạn $10\,\text{m}$; chu kì là $0,5\,\text{s}$.
\choiceTF
{\True Có 5 bước sóng trên đoạn đã cho.}
{\True Bước sóng bằng $2\,\text{m}$.}
{\True Tần số bằng $2\,\text{Hz}$.}
{Tốc độ truyền sóng bằng $1\,\text{m}$ /s.}
\loigiai{
\begin{itemchoice}
\item (Đúng) Có 5 bước sóng trên đoạn đã cho. (Vì): Có 5 bước sóng trên đoạn đã cho. b) Đ: Bước sóng bằng $2\,\text{m}$. c) Đ: Tần số bằng $2\,\text{Hz}$. d) S: Tốc độ truyền sóng bằng $1\,\text{m}$ /s
\item (Đúng) Bước sóng bằng $2\,\text{m}$.
\item (Đúng) Tần số bằng $2\,\text{Hz}$.
\item (Sai) Tốc độ truyền sóng bằng $1\,\text{m}$ /s.
\end{itemchoice}
}
\end{ex}

% ===== Câu 6 =====
% ID: L11C2B8-02-TL
% Mức: TH
\begin{ex}
% ID: L11C2B8-02-TL
% Mức: TH
Trong chuyên đề “Bài toán tổng hợp về các đại lượng đặc trưng của sóng”, Khoảng cách giữa đỉnh thứ nhất và đỉnh thứ mười một là $25\,\text{m}$. Sóng truyền với tốc độ $10\,\text{m}$ /s. Tính bước sóng và tần số.
\choiceTF

\loigiai{
Giữa 11 đỉnh có 10 bước sóng nên $\lambda=25/10=2,5$ m; $f=v/\lambda=4$ Hz.
}
\end{ex}

% ===== Câu 7 =====
% ID: L11C2B8-02-TLN
% Mức: TH
\begin{ex}
% ID: L11C2B8-02-TLN
% Mức: TH
Trong chuyên đề “Bài toán tổng hợp về các đại lượng đặc trưng của sóng”, Khoảng cách giữa 9 đỉnh sóng liên tiếp là $16\,\text{m}$. Bước sóng tính theo m bằng bao nhiêu?
\shortans{2}
\loigiai{
Xác định đúng dữ kiện và đơn vị, áp dụng hệ thức của dạng bài rồi tính được kết quả 2.
}
\end{ex}

% ===== Câu 8 =====
% ID: L11C2B8-02-TN
% Mức: NB
\begin{ex}
% ID: L11C2B8-02-TN
% Mức: NB
Trong chuyên đề “Bài toán tổng hợp về các đại lượng đặc trưng của sóng”, Khoảng cách giữa 7 đỉnh sóng liên tiếp là $12\,\text{m}$. Bước sóng bằng
\choice
{$84\,\text{m}$}
{\True $2\,\text{m}$}
{12/ $7\,\text{m}$}
{$6\,\text{m}$}
\loigiai{
Đáp án B. $2\,\text{m}$. Đối chiếu định nghĩa, hệ thức hoặc dữ kiện của bài toán cho kết luận trên.
}
\end{ex}

% ===== Câu 9 =====
% ID: L11C2B8-03-DS
% Mức: TH
\begin{ex}
% ID: L11C2B8-03-DS
% Mức: TH
Trong chuyên đề “Bài toán tổng hợp về các đại lượng đặc trưng của sóng”, Sóng truyền từ môi trường 1 sang môi trường 2, nguồn không đổi.
\choiceTF
{\True Tần số sóng không đổi.}
{\True Chu kì sóng không đổi.}
{\True Nếu tốc độ giảm thì bước sóng giảm.}
{Tích $v\lambda$ không đổi.}
\loigiai{
\begin{itemchoice}
\item (Đúng) Tần số sóng không đổi. (Vì): Tần số sóng không đổi. b) Đ: Chu kì sóng không đổi. c) Đ: Nếu tốc độ giảm thì bước sóng giảm. d) S: Tích $v\lambda$ không đổi
\item (Đúng) Chu kì sóng không đổi.
\item (Đúng) Nếu tốc độ giảm thì bước sóng giảm.
\item (Sai) Tích $v\lambda$ không đổi.
\end{itemchoice}
}
\end{ex}

% ===== Câu 10 =====
% ID: L11C2B8-03-TL
% Mức: VD
\begin{ex}
% ID: L11C2B8-03-TL
% Mức: VD
Trong chuyên đề “Bài toán tổng hợp về các đại lượng đặc trưng của sóng”, Một sóng đi từ môi trường $A$ sang $B$, tần số $20\,\text{Hz}$ không đổi. Biết tốc độ ở $A$ là $40\,\text{m}$ /s, ở $B$ là $30\,\text{m}$ /s. So sánh bước sóng ở hai môi trường.
\choiceTF

\loigiai{
$\lambda_A=2$ m, $\lambda_B=1,5$ m; bước sóng ở $B$ nhỏ hơn ở $A0,5\,\text{m}$.
}
\end{ex}

% ===== Câu 11 =====
% ID: L11C2B8-03-TLN
% Mức: TH
\begin{ex}
% ID: L11C2B8-03-TLN
% Mức: TH
Trong chuyên đề “Bài toán tổng hợp về các đại lượng đặc trưng của sóng”, Sóng có bước sóng $1,5\,\text{m}$ và chu kì $0,25\,\text{s}$. Tốc độ truyền tính theo m/s bằng bao nhiêu?
\shortans{6}
\loigiai{
Xác định đúng dữ kiện và đơn vị, áp dụng hệ thức của dạng bài rồi tính được kết quả 6.
}
\end{ex}

% ===== Câu 12 =====
% ID: L11C2B8-03-TN
% Mức: TH
\begin{ex}
% ID: L11C2B8-03-TN
% Mức: TH
Trong chuyên đề “Bài toán tổng hợp về các đại lượng đặc trưng của sóng”, Sóng truyền với bước sóng $0,8\,\text{m}$ và chu kì $0,2\,\text{s}$. Tốc độ truyền sóng bằng
\choice
{$1\,\text{m}$ /s}
{$16\,\text{m}$ /s}
{\True $4\,\text{m}$ /s}
{$0,16\,\text{m}$ /s}
\loigiai{
Đáp án C. $4\,\text{m}$ /s. Đối chiếu định nghĩa, hệ thức hoặc dữ kiện của bài toán cho kết luận trên.
}
\end{ex}

% ===== Câu 13 =====
% ID: L11C2B8-04-DS
% Mức: VD
\begin{ex}
% ID: L11C2B8-04-DS
% Mức: VD
Trong chuyên đề “Bài toán tổng hợp về các đại lượng đặc trưng của sóng”, Xét các nhận định tổng hợp thuộc dạng “Bài toán tổng hợp về các đại lượng đặc trưng của sóng”.
\choiceTF
{\True $4\,\text{m}$.}
{\True $2\,\text{m}$.}
{\True $4\,\text{m}$ /s.}
{$100\,\text{m}$.}
\loigiai{
\begin{itemchoice}
\item (Đúng) $4\,\text{m}$. (Vì): $4\,\text{m}$. b) Đ: $2\,\text{m}$. c) Đ: $4\,\text{m}$ /s. d) S: $100\,\text{m}$
\item (Đúng) $2\,\text{m}$.
\item (Đúng) $4\,\text{m}$ /s.
\item (Sai) $100\,\text{m}$.
\end{itemchoice}
}
\end{ex}

% ===== Câu 14 =====
% ID: L11C2B8-04-TL
% Mức: VD
\begin{ex}
% ID: L11C2B8-04-TL
% Mức: VD
Trong chuyên đề “Bài toán tổng hợp về các đại lượng đặc trưng của sóng”, Phân tích các bước lập luận và chỉ ra điều kiện áp dụng trong bài toán: Một sóng có tần số $8\,\text{Hz}$ truyền với tốc độ $24\,\text{m}$ /s. Tính chu kì, bước sóng và quãng đường sóng truyền trong $1,5\,\text{s}$.
\choiceTF

\loigiai{
$T=1/f=0,125$ s; $\lambda=v/f=3$ m; $s=vt=36$ m. Cần nêu rõ điều kiện áp dụng, đổi đơn vị thống nhất và kiểm tra ý nghĩa vật lí của kết quả.
}
\end{ex}

% ===== Câu 15 =====
% ID: L11C2B8-04-TLN
% Mức: VD
\begin{ex}
% ID: L11C2B8-04-TLN
% Mức: VD
Trong chuyên đề “Bài toán tổng hợp về các đại lượng đặc trưng của sóng”, Từ kết quả của tình huống sau, nếu đại lượng cần tìm được nhân với hệ số 2 thì giá trị mới bằng bao nhiêu? Sóng có tốc độ $12\,\text{m}$ /s và tần số $4\,\text{Hz}$. Bước sóng tính theo m bằng bao nhiêu?
\shortans{6}
\loigiai{
Xác định đúng dữ kiện và đơn vị, áp dụng hệ thức của dạng bài rồi tính được kết quả 6.
}
\end{ex}

% ===== Câu 16 =====
% ID: L11C2B8-04-TN
% Mức: TH
\begin{ex}
% ID: L11C2B8-04-TN
% Mức: TH
Trong chuyên đề “Bài toán tổng hợp về các đại lượng đặc trưng của sóng”, Sóng có tốc độ $12\,\text{m}$ /s và tần số $4\,\text{Hz}$. Bước sóng tính theo m bằng bao nhiêu? Chọn giá trị đúng.
\choice
{6}
{01/05/2026}
{4}
{\True 3}
\loigiai{
Đáp án D. 3. Đối chiếu định nghĩa, hệ thức hoặc dữ kiện của bài toán cho kết luận trên.
}
\end{ex}

% ===== Câu 17 =====
% ID: L11C2B8-05-DS
% Mức: VD
\begin{ex}
% ID: L11C2B8-05-DS
% Mức: VD
Trong chuyên đề “Bài toán tổng hợp về các đại lượng đặc trưng của sóng”, Đánh giá cách xử lí các dữ kiện định lượng của dạng “Bài toán tổng hợp về các đại lượng đặc trưng của sóng”.
\choiceTF
{\True Kết quả của bài toán thứ nhất là 3.}
{\True Kết quả của bài toán thứ hai là 2.}
{\True Kết quả của bài toán thứ ba là 6.}
{Có thể bỏ qua hoàn toàn đơn vị và điều kiện vật lí khi tính toán.}
\loigiai{
\begin{itemchoice}
\item (Đúng) Kết quả của bài toán thứ nhất là 3. (Vì): Kết quả của bài toán thứ nhất là 3. b) Đ: Kết quả của bài toán thứ hai là 2. c) Đ: Kết quả của bài toán thứ ba là 6. d) S: Có thể bỏ qua hoàn toàn đơn vị và điều kiện vật lí khi tính toán
\item (Đúng) Kết quả của bài toán thứ hai là 2.
\item (Đúng) Kết quả của bài toán thứ ba là 6.
\item (Sai) Có thể bỏ qua hoàn toàn đơn vị và điều kiện vật lí khi tính toán.
\end{itemchoice}
}
\end{ex}

% ===== Câu 18 =====
% ID: L11C2B8-05-TL
% Mức: VD
\begin{ex}
% ID: L11C2B8-05-TL
% Mức: VD
Trong chuyên đề “Bài toán tổng hợp về các đại lượng đặc trưng của sóng”, Một học sinh giải bài sau nhưng bỏ qua điều kiện vật lí. Hãy sửa cách làm và trình bày lời giải đúng: Khoảng cách giữa đỉnh thứ nhất và đỉnh thứ mười một là $25\,\text{m}$. Sóng truyền với tốc độ $10\,\text{m}$ /s. Tính bước sóng và tần số.
\choiceTF

\loigiai{
Cách giải đúng: Giữa 11 đỉnh có 10 bước sóng nên $\lambda=25/10=2,5$ m; $f=v/\lambda=4$ Hz. Sai lầm cần tránh là dùng hệ thức khi chưa xác định điều kiện biên, môi trường hoặc đơn vị.
}
\end{ex}

% ===== Câu 19 =====
% ID: L11C2B8-05-TLN
% Mức: VD
\begin{ex}
% ID: L11C2B8-05-TLN
% Mức: VD
Trong chuyên đề “Bài toán tổng hợp về các đại lượng đặc trưng của sóng”, Từ kết quả của tình huống sau, nếu đại lượng cần tìm được nhân với hệ số 0.5 thì giá trị mới bằng bao nhiêu? Khoảng cách giữa 9 đỉnh sóng liên tiếp là $16\,\text{m}$. Bước sóng tính theo m bằng bao nhiêu?
\shortans{1}
\loigiai{
Xác định đúng dữ kiện và đơn vị, áp dụng hệ thức của dạng bài rồi tính được kết quả 1.
}
\end{ex}

% ===== Câu 20 =====
% ID: L11C2B8-05-TN
% Mức: TH
\begin{ex}
% ID: L11C2B8-05-TN
% Mức: TH
Trong chuyên đề “Bài toán tổng hợp về các đại lượng đặc trưng của sóng”, Khoảng cách giữa 9 đỉnh sóng liên tiếp là $16\,\text{m}$. Bước sóng tính theo m bằng bao nhiêu? Chọn giá trị đúng.
\choice
{\True 2}
{4}
{1}
{3}
\loigiai{
Đáp án A. 2. Đối chiếu định nghĩa, hệ thức hoặc dữ kiện của bài toán cho kết luận trên.
}
\end{ex}

% ===== Câu 21 =====
% ID: L11C2B8-06-DS
% Mức: TH
\dangbt{Khai thác đồ thị sóng và quan hệ pha giữa các điểm}
\begin{ex}
% ID: L11C2B8-06-DS
% Mức: TH
Hai điểm $M$, $N$ nằm trên cùng phương truyền và cách nhau $3\lambda/4$.
\choiceTF
{\True Độ lệch pha có độ lớn $3\pi/2$ rad.}
{\True $M$ và $N$ vuông pha nếu xét theo môđun $2\pi$.}
{\True Khoảng cách này bằng ba lần khoảng cách gần nhất của hai điểm vuông pha.}
{$M$ và $N$ luôn cùng li độ.}
\loigiai{
\begin{itemchoice}
\item (Đúng) Độ lệch pha có độ lớn $3\pi/2$ rad. (Vì): ộ lệch pha có độ lớn $3\pi/2$ rad. b) Đ: $M$ và $N$ vuông pha nếu xét theo môđun $2\pi$. c) Đ: Khoảng cách này bằng ba lần khoảng cách gần nhất của hai điểm vuông pha. d) S: $M$ và $N$ luôn cùng li độ
\item (Đúng) $M$ và $N$ vuông pha nếu xét theo môđun $2\pi$.
\item (Đúng) Khoảng cách này bằng ba lần khoảng cách gần nhất của hai điểm vuông pha.
\item (Sai) $M$ và $N$ luôn cùng li độ.
\end{itemchoice}
}
\end{ex}

% ===== Câu 22 =====
% ID: L11C2B8-06-TL
% Mức: VD
\dangbt{Bài toán tổng hợp về các đại lượng đặc trưng của sóng}
\begin{ex}
% ID: L11C2B8-06-TL
% Mức: VD
Trong chuyên đề “Bài toán tổng hợp về các đại lượng đặc trưng của sóng”, Mở rộng tình huống sau thành một ví dụ thực tế và trình bày phương pháp giải: Một sóng đi từ môi trường $A$ sang $B$, tần số $20\,\text{Hz}$ không đổi. Biết tốc độ ở $A$ là $40\,\text{m}$ /s, ở $B$ là $30\,\text{m}$ /s. So sánh bước sóng ở hai môi trường.
\choiceTF

\loigiai{
$\lambda_A=2$ m, $\lambda_B=1,5$ m; bước sóng ở $B$ nhỏ hơn ở $A0,5\,\text{m}$. Có thể kiểm chứng bằng cách thay kết quả vào dữ kiện ban đầu và đối chiếu với hiện tượng thực tế.
}
\end{ex}

% ===== Câu 23 =====
% ID: L11C2B8-06-TLN
% Mức: VD
\begin{ex}
% ID: L11C2B8-06-TLN
% Mức: VD
Trong chuyên đề “Bài toán tổng hợp về các đại lượng đặc trưng của sóng”, Từ kết quả của tình huống sau, nếu đại lượng cần tìm được nhân với hệ số 1.5 thì giá trị mới bằng bao nhiêu? Sóng có bước sóng $1,5\,\text{m}$ và chu kì $0,25\,\text{s}$. Tốc độ truyền tính theo m/s bằng bao nhiêu?
\shortans{9}
\loigiai{
Xác định đúng dữ kiện và đơn vị, áp dụng hệ thức của dạng bài rồi tính được kết quả 9.
}
\end{ex}

% ===== Câu 24 =====
% ID: L11C2B8-06-TN
% Mức: TH
\begin{ex}
% ID: L11C2B8-06-TN
% Mức: TH
Trong chuyên đề “Bài toán tổng hợp về các đại lượng đặc trưng của sóng”, Sóng có bước sóng $1,5\,\text{m}$ và chu kì $0,25\,\text{s}$. Tốc độ truyền tính theo m/s bằng bao nhiêu? Chọn giá trị đúng.
\choice
{7}
{\True 6}
{12}
{3}
\loigiai{
Đáp án B. 6. Đối chiếu định nghĩa, hệ thức hoặc dữ kiện của bài toán cho kết luận trên.
}
\end{ex}

% ===== Câu 25 =====
% ID: L11C2B8-07-DS
% Mức: TH
\dangbt{Khai thác đồ thị sóng và quan hệ pha giữa các điểm}
\begin{ex}
% ID: L11C2B8-07-DS
% Mức: TH
Ảnh chụp sóng cho thấy hai đỉnh liên tiếp cách nhau $40\,\text{cm}$.
\choiceTF
{\True Bước sóng bằng $40\,\text{cm}$.}
{\True Đỉnh và đáy gần nhất cách nhau $20\,\text{cm}$.}
{\True Hai điểm cách nhau $10\,\text{cm}$ vuông pha.}
{Hai điểm cách nhau $30\,\text{cm}$ cùng pha.}
\loigiai{
\begin{itemchoice}
\item (Đúng) Bước sóng bằng $40\,\text{cm}$. (Vì): Bước sóng bằng $40\,\text{cm}$. b) Đ: Đỉnh và đáy gần nhất cách nhau $20\,\text{cm}$. c) Đ: Hai điểm cách nhau $10\,\text{cm}$ vuông pha. d) S: Hai điểm cách nhau $30\,\text{cm}$ cùng pha
\item (Đúng) Đỉnh và đáy gần nhất cách nhau $20\,\text{cm}$.
\item (Đúng) Hai điểm cách nhau $10\,\text{cm}$ vuông pha.
\item (Sai) Hai điểm cách nhau $30\,\text{cm}$ cùng pha.
\end{itemchoice}
}
\end{ex}

% ===== Câu 26 =====
% ID: L11C2B8-07-TL
% Mức: TH
\dangbt{Chưa có dạng}
\begin{ex}
% ID: L11C2B8-07-TL
% Mức: TH
Một mũi nhọn $S$ chạm nhẹ vào mặt nước dao động điều hoà với tần số $f =$ 40 $\mathrm{Hz}$. Người ta thấy rằng hai điểm $A$ và $B$ trên mặt nước cùng nằm trên phương truyền sóng cách nhau một khoảng $d =$ 20 $\mathrm{cm}$ luôn dao động ngược pha nhau $. Biết tốc độ truyền sóng nằm trong khoảng từ$ 3 \mathrm{m/s} $đến$ 5 \mathrm{m/s} $. Xác định tốc độ truyền sóng.$ \nguon{ly11kntt.pdf}
\choiceTF

\loigiai{
Vì sóng tại hai điểm $A$, $B$ ngược pha nhau, nên khoảng cách AB thoả mãn:

AB = d = (2k+1) $\dfrac{\lambda}{2}$ = (2k+1) $\dfrac{v}{$ 2 $\mathrm{f}$ }\Rightarrow $v =$ \dfrac{2fd}{2k + 1} = \dfrac{16}{2k + 1} $, với k∈$ Z $.

Theo đề bài:$3\,\mathrm{m}$/s≤v≤$5\,\mathrm{m}$/s \Rightarrow  3≤$\dfrac{16}{2k + 1}$≤5 ⇔1,1≤k≤2,17

Vậy k=2 . Suy ra tốc độ truyền sóng là:$ v = \dfrac{16}{2k + 1} = \dfrac{16}{2 \cdot 2 + 1} = \dfrac{16}{5} = 3,2\text{m/s}.
}
\end{ex}

% ===== Câu 27 =====
% ID: L11C2B8-07-TLN
% Mức: TH
\dangbt{Khai thác đồ thị sóng và quan hệ pha giữa các điểm}
\begin{ex}
% ID: L11C2B8-07-TLN
% Mức: TH
Hai điểm cách nhau $0,6\,\text{m}$ có độ lệch pha $\pi$ rad và là hai điểm ngược pha gần nhau nhất. Bước sóng tính theo m bằng bao nhiêu?
\shortans{1,2}
\loigiai{
Xác định đúng dữ kiện và đơn vị, áp dụng hệ thức của dạng bài rồi tính được kết quả 1.2.
}
\end{ex}

% ===== Câu 28 =====
% ID: L11C2B8-07-TN
% Mức: VD
\dangbt{Bài toán tổng hợp về các đại lượng đặc trưng của sóng}
\begin{ex}
% ID: L11C2B8-07-TN
% Mức: VD
Trong chuyên đề “Bài toán tổng hợp về các đại lượng đặc trưng của sóng”, Kết luận nào phù hợp nhất khi giải bài toán sau: Một sóng có tần số $8\,\text{Hz}$ truyền với tốc độ $24\,\text{m}$ /s. Tính chu kì, bước sóng và quãng đường sóng truyền trong $1,5\,\text{s}$.
\choice
{Đại lượng cần tìm luôn bằng 0 trong mọi trường hợp.}
{Kết quả không phụ thuộc môi trường, tần số hoặc điều kiện biên.}
{\True $T=1/f=0,125$ s}
{Không cần sử dụng định nghĩa hay dữ kiện đã cho.}
\loigiai{
Đáp án C. $T=1/f=0,125$ s. Đối chiếu định nghĩa, hệ thức hoặc dữ kiện của bài toán cho kết luận trên.
}
\end{ex}

% ===== Câu 29 =====
% ID: L11C2B8-08-DS
% Mức: TH
\dangbt{Khai thác đồ thị sóng và quan hệ pha giữa các điểm}
\begin{ex}
% ID: L11C2B8-08-DS
% Mức: TH
Sóng truyền theo chiều dương Ox.
\choiceTF
{\True Điểm gần nguồn hơn dao động sớm pha hơn điểm xa nguồn.}
{\True Pha giảm khi tọa độ tăng tại cùng thời điểm.}
{\True Hai điểm cách nhau một bước sóng có cùng trạng thái dao động.}
{Mọi điểm trên phương truyền đều có cùng pha.}
\loigiai{
\begin{itemchoice}
\item (Đúng) Điểm gần nguồn hơn dao động sớm pha hơn điểm xa nguồn. (Vì): iểm gần nguồn hơn dao động sớm pha hơn điểm xa nguồn. b) Đ: Pha giảm khi tọa độ tăng tại cùng thời điểm. c) Đ: Hai điểm cách nhau một bước sóng có cùng trạng thái dao động. d) S: Mọi điểm trên phương truyền đều có cùng pha
\item (Đúng) Pha giảm khi tọa độ tăng tại cùng thời điểm.
\item (Đúng) Hai điểm cách nhau một bước sóng có cùng trạng thái dao động.
\item (Sai) Mọi điểm trên phương truyền đều có cùng pha.
\end{itemchoice}
}
\end{ex}

% ===== Câu 30 =====
% ID: L11C2B8-08-TL
% Mức: TH
\dangbt{Chưa có dạng}
\begin{ex}
% ID: L11C2B8-08-TL
% Mức: TH
Trong môi trường đàn hồi, có một sóng cơ tần số $10 \mathrm{Hz}$ lan truyền với tốc độ $40 \mathrm{cm}$ /s. Hai điểm $A$, $B$ trên phương truyền sóng dao động cùng pha nhau. Giữa chúng chỉ có hai điểm khác dao động ngược pha với A. Tính khoảng cách AB. $\nguon{ly11 kntt.pdf}$
\choiceTF

\loigiai{
Hai điểm $A$, $B$ dao động cùng pha, nên: AB = λ;2λ;3λ;... nhưng giữa chúng chỉ có hai điểm ngược pha với $A$ nên: AB = 2λ = $\dfrac{$ 2 $\mathrm{v}$ }{f} = \dfrac{2.40}{10} $=$ 8\,\mathrm{cm}.
}
\end{ex}

% ===== Câu 31 =====
% ID: L11C2B8-08-TLN
% Mức: TH
\dangbt{Khai thác đồ thị sóng và quan hệ pha giữa các điểm}
\begin{ex}
% ID: L11C2B8-08-TLN
% Mức: TH
Bước sóng là $2,4\,\text{m}$. Khoảng cách gần nhất giữa hai điểm vuông pha tính theo m bằng bao nhiêu?
\shortans{0.6}
\loigiai{
Xác định đúng dữ kiện và đơn vị, áp dụng hệ thức của dạng bài rồi tính được kết quả 0.6.
}
\end{ex}

% ===== Câu 32 =====
% ID: L11C2B8-08-TN
% Mức: VD
\dangbt{Bài toán tổng hợp về các đại lượng đặc trưng của sóng}
\begin{ex}
% ID: L11C2B8-08-TN
% Mức: VD
Trong chuyên đề “Bài toán tổng hợp về các đại lượng đặc trưng của sóng”, Kết luận nào phù hợp nhất khi giải bài toán sau: Khoảng cách giữa đỉnh thứ nhất và đỉnh thứ mười một là $25\,\text{m}$. Sóng truyền với tốc độ $10\,\text{m}$ /s. Tính bước sóng và tần số.
\choice
{Không cần sử dụng định nghĩa hay dữ kiện đã cho.}
{Đại lượng cần tìm luôn bằng 0 trong mọi trường hợp.}
{Kết quả không phụ thuộc môi trường, tần số hoặc điều kiện biên.}
{\True Giữa 11 đỉnh có 10 bước sóng nên $\lambda=25/10=2,5$ m}
\loigiai{
Đáp án D. Giữa 11 đỉnh có 10 bước sóng nên $\lambda=25/10=2,5$ m. Đối chiếu định nghĩa, hệ thức hoặc dữ kiện của bài toán cho kết luận trên.
}
\end{ex}

% ===== Câu 33 =====
% ID: L11C2B8-09-DS
% Mức: VD
\dangbt{Khai thác đồ thị sóng và quan hệ pha giữa các điểm}
\begin{ex}
% ID: L11C2B8-09-DS
% Mức: VD
Xét các nhận định tổng hợp thuộc dạng “Khai thác đồ thị sóng và quan hệ pha giữa các điểm”.
\choiceTF
{\True một bước sóng.}
{\True ngược pha.}
{\True $\lambda/2$.}
{nửa bước sóng.}
\loigiai{
\begin{itemchoice}
\item (Đúng) một bước sóng. (Vì): một bước sóng. b) Đ: ngược pha. c) Đ: $\lambda/2$. d) S: nửa bước sóng
\item (Đúng) ngược pha.
\item (Đúng) $\lambda/2$.
\item (Sai) nửa bước sóng.
\end{itemchoice}
}
\end{ex}

% ===== Câu 34 =====
% ID: L11C2B8-09-TL
% Mức: TH
\dangbt{Chưa có dạng}
\begin{ex}
% ID: L11C2B8-09-TL
% Mức: TH
Trong môi trường đàn hồi, có một sóng cơ có tần số $10 \mathrm{Hz}$ lan truyền với tốc độ $40 \mathrm{cm}$ /s. Hai điểm $A$, $B$ trên phương truyền sóng dao động cùng pha nhau. Giữa chúng có hai điểm $M$ và N. Biết rằng khi $M$ hoặc $N$ có tốc độ dao động cực đại thì tại $A$ tốc độ dao động cực tiểu. Tính khoảng cách AB. $\nguon{ly11 kntt.pdf}$
\choiceTF

\loigiai{
Theo đề bài, khi $M$ hoặc $N$ có tốc độ dao động cực đại thì tại $A$ có tốc độ dao động cực tiểu, tức là $M$, $N$ dao động vuông pha với A.

Hai điểm $A$, $B$ dao động cùng pha, nên: AB = λ;2λ;3λ;... nhưng giữa chúng chỉ có hai điểm dao động vuông pha với $A$ nên: AB = λ = $\dfrac{v}{f}$ = $\dfrac{40}{10}$ = $4\,\mathrm{cm}$.
}
\end{ex}

% ===== Câu 35 =====
% ID: L11C2B8-09-TLN
% Mức: TH
\dangbt{Khai thác đồ thị sóng và quan hệ pha giữa các điểm}
\begin{ex}
% ID: L11C2B8-09-TLN
% Mức: TH
Hai đỉnh liên tiếp trên ảnh sóng cách nhau $75\,\text{cm}$. Khoảng cách từ một đỉnh đến đáy gần nhất tính theo cm bằng bao nhiêu?
\shortans{37.5}
\loigiai{
Xác định đúng dữ kiện và đơn vị, áp dụng hệ thức của dạng bài rồi tính được kết quả 37.5.
}
\end{ex}

% ===== Câu 36 =====
% ID: L11C2B8-09-TN
% Mức: VD
\dangbt{Bài toán tổng hợp về các đại lượng đặc trưng của sóng}
\begin{ex}
% ID: L11C2B8-09-TN
% Mức: VD
Trong chuyên đề “Bài toán tổng hợp về các đại lượng đặc trưng của sóng”, Kết luận nào phù hợp nhất khi giải bài toán sau: Một sóng đi từ môi trường $A$ sang $B$, tần số $20\,\text{Hz}$ không đổi. Biết tốc độ ở $A$ là $40\,\text{m}$ /s, ở $B$ là $30\,\text{m}$ /s. So sánh bước sóng ở hai môi trường.
\choice
{\True $\lambda_A=2$ m, $\lambda_B=1,5$ m}
{Không cần sử dụng định nghĩa hay dữ kiện đã cho.}
{Đại lượng cần tìm luôn bằng 0 trong mọi trường hợp.}
{Kết quả không phụ thuộc môi trường, tần số hoặc điều kiện biên.}
\loigiai{
Đáp án A. $\lambda_A=2$ m, $\lambda_B=1,5$ m. Đối chiếu định nghĩa, hệ thức hoặc dữ kiện của bài toán cho kết luận trên.
}
\end{ex}

% ===== Câu 37 =====
% ID: L11C2B8-10-DS
% Mức: VD
\dangbt{Khai thác đồ thị sóng và quan hệ pha giữa các điểm}
\begin{ex}
% ID: L11C2B8-10-DS
% Mức: VD
Đánh giá cách xử lí các dữ kiện định lượng của dạng “Khai thác đồ thị sóng và quan hệ pha giữa các điểm”.
\choiceTF
{\True Kết quả của bài toán thứ nhất là 1.2.}
{\True Kết quả của bài toán thứ hai là 0.6.}
{\True Kết quả của bài toán thứ ba là 37.5.}
{Có thể bỏ qua hoàn toàn đơn vị và điều kiện vật lí khi tính toán.}
\loigiai{
\begin{itemchoice}
\item (Đúng) Kết quả của bài toán thứ nhất là 1.2. (Vì): Kết quả của bài toán thứ nhất là 1.2. b) Đ: Kết quả của bài toán thứ hai là 0.6. c) Đ: Kết quả của bài toán thứ ba là 37.5. d) S: Có thể bỏ qua hoàn toàn đơn vị và điều kiện vật lí khi tính toán
\item (Đúng) Kết quả của bài toán thứ hai là 0.6.
\item (Đúng) Kết quả của bài toán thứ ba là 37.5.
\item (Sai) Có thể bỏ qua hoàn toàn đơn vị và điều kiện vật lí khi tính toán.
\end{itemchoice}
}
\end{ex}

% ===== Câu 38 =====
% ID: L11C2B8-10-TL
% Mức: TH
\dangbt{Chưa có dạng}
\begin{ex}
% ID: L11C2B8-10-TL
% Mức: TH
*. Một sóng cơ lan truyền qua điểm $M$ rồi đến điểm $N$ cùng nằm trên một phương truyền sóng cách nhau một phần ba bước sóng. Tại thời điểm $t = 0 li độ$ tại $M$ là + $4 \mathrm{cm}$ và tại $N$ là $-4 \mathrm{cm}$. Xác định thời điểm t_1 và t $2\,\mathrm{g}$ ần nhất để $M$ và $N$ lên đến vị trí cao nhất. Biết chu kì sóng là $T = 1\$, $\mathrm{s}$. $\nguon{ly11 kntt.pdf}$
\choiceTF

\loigiai{
Sử dụng đồ thị li độ - quãng đường Hình 8.1 $G$ của sóng quy ước chiều truyền dương để xác định các vùng mà các phần tử vật chất đang đi lên và đi xuống. Li độ

[Một sóng cơ lan truyền qua điểm $M$ rồi đến điểm N]

Vì sóng truyền qua $M$ rồi mới đến $N$, nên $M$ ở bên trái và $N$ ở bên phải, mặt khác vì u_(M) = + $4\,\mathrm{cm}$ và u_(N) = - $4\,\mathrm{cm}$, nên chúng phải nằm ở vị trí như Hình 8.1 $G$ (cả $M$ và $N$ đều đang đi lên).

Vì $M$ cách đỉnh gần nhất một khoảng là $\dfrac{\lambda}{12}$ nên thời gian ngắn nhất để $M$ đi từ vị trí hiện tại đến vị trí cao nhất là t₁ = $\dfrac{T}{12} = \dfrac{1}{12}$ s.

Tương tự ta xác định được, thời gian ngắn nhất để $N$ đến vị trí cân bằng là $\dfrac{T}{6}$ và thời gian ngắn nhất để đi từ vị trí cân bằng đến vị trí cao nhất là $\dfrac{T}{4}$ nên t₂ = $\dfrac{\text{T}}{6} + \dfrac{\text{T}}{4} = \dfrac{5\text{T}}{12} = \dfrac{5}{12}\text{s}$.
}
\end{ex}

% ===== Câu 39 =====
% ID: L11C2B8-10-TLN
% Mức: VD
\dangbt{Khai thác đồ thị sóng và quan hệ pha giữa các điểm}
\begin{ex}
% ID: L11C2B8-10-TLN
% Mức: VD
Từ kết quả của tình huống sau, nếu đại lượng cần tìm được nhân với hệ số 2 thì giá trị mới bằng bao nhiêu? Hai điểm cách nhau $0,6\,\text{m}$ có độ lệch pha $\pi$ rad và là hai điểm ngược pha gần nhau nhất. Bước sóng tính theo m bằng bao nhiêu?
\shortans{02/04/2026}
\loigiai{
Xác định đúng dữ kiện và đơn vị, áp dụng hệ thức của dạng bài rồi tính được kết quả 2.4.
}
\end{ex}

% ===== Câu 40 =====
% ID: L11C2B8-10-TN
% Mức: VD
\dangbt{Bài toán tổng hợp về các đại lượng đặc trưng của sóng}
\begin{ex}
% ID: L11C2B8-10-TN
% Mức: VD
Trong chuyên đề “Bài toán tổng hợp về các đại lượng đặc trưng của sóng”, Phương pháp phù hợp nhất cho dạng “Bài toán tổng hợp về các đại lượng đặc trưng của sóng” là
\choice
{luôn coi mọi điểm dao động cùng pha}
{\True xác định dữ kiện, chọn đúng hệ thức hoặc định nghĩa, đổi đơn vị rồi kiểm tra kết quả}
{chỉ thay số mà không cần xét điều kiện áp dụng}
{bỏ qua đơn vị và chiều truyền sóng}
\loigiai{
Đáp án B. xác định dữ kiện, chọn đúng hệ thức hoặc định nghĩa, đổi đơn vị rồi kiểm tra kết quả. Đối chiếu định nghĩa, hệ thức hoặc dữ kiện của bài toán cho kết luận trên.
}
\end{ex}

% ===== Câu 41 =====
% ID: L11C2B8-11-DS
% Mức: TH
\dangbt{Nhận biết bản chất và sự truyền sóng cơ}
\begin{ex}
% ID: L11C2B8-11-DS
% Mức: TH
Xét một sóng cơ truyền trên mặt nước.
\choiceTF
{\True Nguồn sóng cung cấp dao động cho môi trường.}
{\True Các phần tử nước dao động quanh vị trí cân bằng.}
{\True Sóng mang năng lượng ra xa nguồn.}
{Các phần tử nước chuyển động cùng sóng đến bờ.}
\loigiai{
\begin{itemchoice}
\item (Đúng) Nguồn sóng cung cấp dao động cho môi trường. (Vì): Nguồn sóng cung cấp dao động cho môi trường. b) Đ: Các phần tử nước dao động quanh vị trí cân bằng. c) Đ: Sóng mang năng lượng ra xa nguồn. d) S: Các phần tử nước chuyển động cùng sóng đến bờ
\item (Đúng) Các phần tử nước dao động quanh vị trí cân bằng.
\item (Đúng) Sóng mang năng lượng ra xa nguồn.
\item (Sai) Các phần tử nước chuyển động cùng sóng đến bờ.
\end{itemchoice}
}
\end{ex}

% ===== Câu 42 =====
% ID: L11C2B8-11-TL
% Mức: TH
\dangbt{Chưa có dạng}
\begin{ex}
% ID: L11C2B8-11-TL
% Mức: TH
*. Trên mặt thoáng của một chất lỏng, một mũi nhọn $O$ chạm vào mặt thoáng dao động điều hoà với tần số f, tạo thành sóng trên mặt thoáng với bước sóng $\lambda$. Xét hai phương truyền sóng Ox và O𝑦 vuông góc với nhau. Gọi $M$ là một điểm thuộc Ox cách $O$ một đoạn 16 $\lambda$ và $N$ thuộc Oy cách $O$ một đoạn 12 $\lambda$. Tính số điểm dao động đồng pha với nguồn $O$ trên đoạn MN (không kể $M$, N). BÀl 9. SÓNG NGANG. SÓNG DỌC. SỰ TRUYỀ $N$ NĂNG LƯỢNG CỦ $A$ SÓNG CƠ $\nguon{ly11 kntt.pdf}$
\choiceTF

\loigiai{
Vị trí của các điểm $O$, $M$, $N$ được mô tả như trên Hình $8.2\,\text{G}$. Kẻ OH⊥MN, $\Delta$ OMN vuông nên ta có: $\left. \dfrac{1}{\text{OH}^{2}} = \dfrac{1}{\text{OM}^{2}} + \dfrac{1}{\text{ON}^{2}}\Rightarrow\text{OH} = 9,6\lambda \right.$

[Trên mặt thoáng của một chất lỏng, một mũi nhọn O]

Các điểm dao động cùng pha với $O$, cách $O$ những khoảng: d = kλ .

Xét trên đoạn MH có: 9,6λ≤kλ≤16λ $\Rightarrow$ 9,6≤k≤16

$\Rightarrow$  k = 10,11,...,16, vậy trên MH có 7 điểm.

Xét trên đoạn NH có: 9,6λ≤kλ≤12λ$\Rightarrow$ 9,6≤k≤12

$\Rightarrow$ k = 10,11,12, vậy trên MH có 3 điểm.

Như vậy, tổng số điểm dao động cùng pha với $O$ trên MN là 10 điểm.
}
\end{ex}

% ===== Câu 43 =====
% ID: L11C2B8-11-TLN
% Mức: VD
\dangbt{Khai thác đồ thị sóng và quan hệ pha giữa các điểm}
\begin{ex}
% ID: L11C2B8-11-TLN
% Mức: VD
Từ kết quả của tình huống sau, nếu đại lượng cần tìm được nhân với hệ số 0.5 thì giá trị mới bằng bao nhiêu? Bước sóng là $2,4\,\text{m}$. Khoảng cách gần nhất giữa hai điểm vuông pha tính theo m bằng bao nhiêu?
\shortans{0.3}
\loigiai{
Xác định đúng dữ kiện và đơn vị, áp dụng hệ thức của dạng bài rồi tính được kết quả 0.3.
}
\end{ex}

% ===== Câu 44 =====
% ID: L11C2B8-11-TN
% Mức: NB
\dangbt{Chưa có dạng}
\begin{ex}
% ID: L11C2B8-11-TN
% Mức: NB
Tại một điểm $O$ trên mặt nước có một nguồn dao động điều hoà theo phương thẳng đứng với tần số $2 \mathrm{Hz}$. Từ điểm $O$ có những gợn sóng tròn lan rộng ra xung quanh. Khoảng cách giữa hai gợn sóng kế tiếp là $20 \mathrm{cm}$. Tốc độ truyền sóng trên mặt nước là $\nguon{ly11 kntt.pdf}$
\choice
{$20 \mathrm{cm}$ /s.}
{\True $40 \mathrm{cm}$ /s.}
{$80 \mathrm{cm}$ /s.}
{$120\,\text{cm}$ /s}
\loigiai{
Chọn B.

Khoảng cách giữa hai gợn sóng bằng 1 bước sóng nên λ = $20\,\mathrm{cm}$ Tốc độ truyền sóng: v = λf = 20.2 = $40\,\mathrm{cm}$ /s
}
\end{ex}

% ===== Câu 45 =====
% ID: L11C2B8-12-DS
% Mức: TH
\dangbt{Nhận biết bản chất và sự truyền sóng cơ}
\begin{ex}
% ID: L11C2B8-12-DS
% Mức: TH
Một loa phát âm trong phòng kín.
\choiceTF
{\True Âm là sóng cơ.}
{\True Không khí là môi trường truyền âm.}
{Âm truyền được trong chân không.}
{\True Khi âm tới tai, màng nhĩ nhận năng lượng.}
\loigiai{
\begin{itemchoice}
\item (Đúng) Âm là sóng cơ. (Vì): Âm là sóng cơ. b) Đ: Không khí là môi trường truyền âm. c) S: Âm truyền được trong chân không. d) Đ: Khi âm tới tai, màng nhĩ nhận năng lượng
\item (Đúng) Không khí là môi trường truyền âm.
\item (Sai) Âm truyền được trong chân không.
\item (Đúng) Khi âm tới tai, màng nhĩ nhận năng lượng.
\end{itemchoice}
}
\end{ex}

% ===== Câu 46 =====
% ID: L11C2B8-12-TL
% Mức: TH
\dangbt{Khai thác đồ thị sóng và quan hệ pha giữa các điểm}
\begin{ex}
% ID: L11C2B8-12-TL
% Mức: TH
Chứng minh hai điểm cách nhau một số nguyên lần bước sóng dao động cùng pha.
\choiceTF

\loigiai{
Vì $\Delta\varphi=2\pi d/\lambda$. Với $d=k\lambda$ thì $\Delta\varphi=2k\pi$, nên hai điểm cùng pha.
}
\end{ex}

% ===== Câu 47 =====
% ID: L11C2B8-12-TLN
% Mức: VD
\begin{ex}
% ID: L11C2B8-12-TLN
% Mức: VD
Từ kết quả của tình huống sau, nếu đại lượng cần tìm được nhân với hệ số 1.5 thì giá trị mới bằng bao nhiêu? Hai đỉnh liên tiếp trên ảnh sóng cách nhau $75\,\text{cm}$. Khoảng cách từ một đỉnh đến đáy gần nhất tính theo cm bằng bao nhiêu?
\shortans{56.25}
\loigiai{
Xác định đúng dữ kiện và đơn vị, áp dụng hệ thức của dạng bài rồi tính được kết quả 56.25.
}
\end{ex}

% ===== Câu 48 =====
% ID: L11C2B8-12-TN
% Mức: NB
\dangbt{Chưa có dạng}
\begin{ex}
% ID: L11C2B8-12-TN
% Mức: NB
Một sóng có tần số $120 \mathrm{Hz}$ truyền trong một môi trường với tốc độ $60 \mathrm{m/s}$. Bước sóng của nó là $\nguon{ly11 kntt.pdf}$
\choice
{$1,0\,\text{m}$.}
{$2,0\,\text{m}$.}
{\True $0,5\,\text{m}$.}
{$0,25\,\text{m}$.}
\loigiai{
Chọn C.

Bước sóng λ = $\dfrac{v}{f}=\dfrac{60}{120}=0,5\,\mathrm{m}$
}
\end{ex}

% ===== Câu 49 =====
% ID: L11C2B8-13-DS
% Mức: TH
\dangbt{Nhận biết bản chất và sự truyền sóng cơ}
\begin{ex}
% ID: L11C2B8-13-DS
% Mức: TH
Quan sát một xung truyền trên dây căng.
\choiceTF
{\True Biến dạng lan dọc theo dây.}
{\True Mỗi phần tử dây chỉ dao động quanh vị trí cân bằng.}
{Tốc độ truyền xung luôn bằng tốc độ dao động của phần tử dây.}
{\True Lực liên kết giữa các phần tử giúp truyền dao động.}
\loigiai{
\begin{itemchoice}
\item (Đúng) Biến dạng lan dọc theo dây. (Vì): Biến dạng lan dọc theo dây. b) Đ: Mỗi phần tử dây chỉ dao động quanh vị trí cân bằng. c) S: Tốc độ truyền xung luôn bằng tốc độ dao động của phần tử dây. d) Đ: Lực liên kết giữa các phần tử giúp truyền dao động
\item (Đúng) Mỗi phần tử dây chỉ dao động quanh vị trí cân bằng.
\item (Sai) Tốc độ truyền xung luôn bằng tốc độ dao động của phần tử dây.
\item (Đúng) Lực liên kết giữa các phần tử giúp truyền dao động.
\end{itemchoice}
}
\end{ex}

% ===== Câu 50 =====
% ID: L11C2B8-13-TL
% Mức: TH
\dangbt{Khai thác đồ thị sóng và quan hệ pha giữa các điểm}
\begin{ex}
% ID: L11C2B8-13-TL
% Mức: TH
Một ảnh sóng có bước sóng $80\,\text{cm}$. Xác định ba vị trí gần nhất tính từ $O$ theo chiều truyền mà phần tử ngược pha với O.
\choiceTF

\loigiai{
Điểm ngược pha thỏa $d=(k+1/2)\lambda$. Ba vị trí gần nhất là $40\,\text{cm}$, $120\,\text{cm}$ và $200\,\text{cm}$.
}
\end{ex}

% ===== Câu 51 =====
% ID: L11C2B8-13-TLN
% Mức: TH
\dangbt{Nhận biết bản chất và sự truyền sóng cơ}
\begin{ex}
% ID: L11C2B8-13-TLN
% Mức: TH
Một sóng truyền quãng đường $18\,\text{m}$ trong $6\,\text{s}$. Tốc độ truyền sóng tính theo m/s bằng bao nhiêu?
\shortans{3}
\loigiai{
Xác định đúng dữ kiện và đơn vị, áp dụng hệ thức của dạng bài rồi tính được kết quả 3.
}
\end{ex}

% ===== Câu 52 =====
% ID: L11C2B8-13-TN
% Mức: NB
\dangbt{Chưa có dạng}
\begin{ex}
% ID: L11C2B8-13-TN
% Mức: NB
Thời gian kể từ khi ngọn sóng thứ nhất đến ngọn sóng thứ sáu đi qua trước mặt một người quan sát là $12\,\mathrm{s}$. Tốc độ truyền sóng là $2 \mathrm{m/s}$. Bước sóng có giá trị là $\nguon{ly11 kntt.pdf}$
\choice
{\True $4,8\,\text{m}$.}
{$4\,\text{m}$.}
{$6 \mathrm{cm}$.}
{$48 \mathrm{cm}$.}
\loigiai{
Chọn A.

Khoảng thời gian từ ngọn thứ nhất đến ngọn thứ sáu ứng với 5 chu kì.

Suy ra 5 $T=12\,\mathrm{s}$ nên $T$ = $2,4\,\mathrm{s}$.

Bước sóng λ = vT = 2.2,4 = $4,8\,\mathrm{m}$
}
\end{ex}

% ===== Câu 53 =====
% ID: L11C2B8-14-DS
% Mức: VD
\dangbt{Nhận biết bản chất và sự truyền sóng cơ}
\begin{ex}
% ID: L11C2B8-14-DS
% Mức: VD
Xét các nhận định tổng hợp thuộc dạng “Nhận biết bản chất và sự truyền sóng cơ”.
\choiceTF
{\True năng lượng và trạng thái dao động.}
{\True chân không.}
{\True dao động quanh vị trí cân bằng.}
{các phần tử vật chất.}
\loigiai{
\begin{itemchoice}
\item (Đúng) năng lượng và trạng thái dao động. (Vì): năng lượng và trạng thái dao động. b) Đ: chân không. c) Đ: dao động quanh vị trí cân bằng. d) S: các phần tử vật chất
\item (Đúng) chân không.
\item (Đúng) dao động quanh vị trí cân bằng.
\item (Sai) các phần tử vật chất.
\end{itemchoice}
}
\end{ex}

% ===== Câu 54 =====
% ID: L11C2B8-14-TL
% Mức: VD
\dangbt{Khai thác đồ thị sóng và quan hệ pha giữa các điểm}
\begin{ex}
% ID: L11C2B8-14-TL
% Mức: VD
Giải thích vì sao đường cong hình sin tại một thời điểm không phải là quỹ đạo chuyển động của một phần tử môi trường.
\choiceTF

\loigiai{
Đường cong biểu diễn li độ của nhiều phần tử ở các vị trí khác nhau tại cùng thời điểm; mỗi phần tử chỉ dao động quanh vị trí cân bằng riêng.
}
\end{ex}

% ===== Câu 55 =====
% ID: L11C2B8-14-TLN
% Mức: TH
\dangbt{Nhận biết bản chất và sự truyền sóng cơ}
\begin{ex}
% ID: L11C2B8-14-TLN
% Mức: TH
Sóng truyền $24\,\text{m}$ trong $8\,\text{s}$. Tốc độ truyền sóng tính theo m/s bằng bao nhiêu?
\shortans{3}
\loigiai{
Xác định đúng dữ kiện và đơn vị, áp dụng hệ thức của dạng bài rồi tính được kết quả 3.
}
\end{ex}

% ===== Câu 56 =====
% ID: L11C2B8-14-TN
% Mức: NB
\dangbt{Khai thác đồ thị sóng và quan hệ pha giữa các điểm}
\begin{ex}
% ID: L11C2B8-14-TN
% Mức: NB
Hai điểm gần nhau nhất trên cùng phương truyền sóng dao động cùng pha cách nhau
\choice
{một phần tư bước sóng}
{hai bước sóng}
{\True một bước sóng}
{nửa bước sóng}
\loigiai{
Một bước sóng. Đối chiếu định nghĩa, hệ thức hoặc dữ kiện của bài toán cho kết luận trên.
}
\end{ex}

% ===== Câu 57 =====
% ID: L11C2B8-15-DS
% Mức: VD
\dangbt{Nhận biết bản chất và sự truyền sóng cơ}
\begin{ex}
% ID: L11C2B8-15-DS
% Mức: VD
Đánh giá cách xử lí các dữ kiện định lượng của dạng “Nhận biết bản chất và sự truyền sóng cơ”.
\choiceTF
{\True Kết quả của bài toán thứ nhất là 3.}
{\True Kết quả của bài toán thứ hai là 3.}
{\True Kết quả của bài toán thứ ba là 3.}
{Có thể bỏ qua hoàn toàn đơn vị và điều kiện vật lí khi tính toán.}
\loigiai{
\begin{itemchoice}
\item (Đúng) Kết quả của bài toán thứ nhất là 3. (Vì): Kết quả của bài toán thứ nhất là 3. b) Đ: Kết quả của bài toán thứ hai là 3. c) Đ: Kết quả của bài toán thứ ba là 3. d) S: Có thể bỏ qua hoàn toàn đơn vị và điều kiện vật lí khi tính toán
\item (Đúng) Kết quả của bài toán thứ hai là 3.
\item (Đúng) Kết quả của bài toán thứ ba là 3.
\item (Sai) Có thể bỏ qua hoàn toàn đơn vị và điều kiện vật lí khi tính toán.
\end{itemchoice}
}
\end{ex}

% ===== Câu 58 =====
% ID: L11C2B8-15-TL
% Mức: VD
\dangbt{Khai thác đồ thị sóng và quan hệ pha giữa các điểm}
\begin{ex}
% ID: L11C2B8-15-TL
% Mức: VD
Phân tích các bước lập luận và chỉ ra điều kiện áp dụng trong bài toán: Chứng minh hai điểm cách nhau một số nguyên lần bước sóng dao động cùng pha.
\choiceTF

\loigiai{
Vì $\Delta\varphi=2\pi d/\lambda$. Với $d=k\lambda$ thì $\Delta\varphi=2k\pi$, nên hai điểm cùng pha. Cần nêu rõ điều kiện áp dụng, đổi đơn vị thống nhất và kiểm tra ý nghĩa vật lí của kết quả.
}
\end{ex}

% ===== Câu 59 =====
% ID: L11C2B8-15-TLN
% Mức: TH
\dangbt{Nhận biết bản chất và sự truyền sóng cơ}
\begin{ex}
% ID: L11C2B8-15-TLN
% Mức: TH
Một xung sóng đi được $15\,\text{m}$ với tốc độ $5\,\text{m}$ /s. Thời gian truyền tính theo s bằng bao nhiêu?
\shortans{3}
\loigiai{
Xác định đúng dữ kiện và đơn vị, áp dụng hệ thức của dạng bài rồi tính được kết quả 3.
}
\end{ex}

% ===== Câu 60 =====
% ID: L11C2B8-15-TN
% Mức: NB
\dangbt{Khai thác đồ thị sóng và quan hệ pha giữa các điểm}
\begin{ex}
% ID: L11C2B8-15-TN
% Mức: NB
Hai điểm cách nhau $\lambda/2$ trên cùng phương truyền sóng dao động
\choice
{cùng pha}
{vuông pha}
{không xác định}
{\True ngược pha}
\loigiai{
Đáp án D. ngược pha. Đối chiếu định nghĩa, hệ thức hoặc dữ kiện của bài toán cho kết luận trên.
}
\end{ex}

% ===== Câu 61 =====
% ID: L11C2B8-16-DS
% Mức: TH
\dangbt{Phân tích trạng thái và chiều truyền sóng từ dữ kiện quan sát}
\begin{ex}
% ID: L11C2B8-16-DS
% Mức: TH
Trong chuyên đề “Phân tích trạng thái và chiều truyền sóng từ dữ kiện quan sát”, Hai điểm $M$, $N$ nằm trên cùng phương truyền và cách nhau $3\lambda/4$.
\choiceTF
{\True Độ lệch pha có độ lớn $3\pi/2$ rad.}
{\True $M$ và $N$ vuông pha nếu xét theo môđun $2\pi$.}
{\True Khoảng cách này bằng ba lần khoảng cách gần nhất của hai điểm vuông pha.}
{$M$ và $N$ luôn cùng li độ.}
\loigiai{
\begin{itemchoice}
\item (Đúng) Độ lệch pha có độ lớn $3\pi/2$ rad. (Vì): ộ lệch pha có độ lớn $3\pi/2$ rad. b) Đ: $M$ và $N$ vuông pha nếu xét theo môđun $2\pi$. c) Đ: Khoảng cách này bằng ba lần khoảng cách gần nhất của hai điểm vuông pha. d) S: $M$ và $N$ luôn cùng li độ
\item (Đúng) $M$ và $N$ vuông pha nếu xét theo môđun $2\pi$.
\item (Đúng) Khoảng cách này bằng ba lần khoảng cách gần nhất của hai điểm vuông pha.
\item (Sai) $M$ và $N$ luôn cùng li độ.
\end{itemchoice}
}
\end{ex}

% ===== Câu 62 =====
% ID: L11C2B8-16-TL
% Mức: VD
\dangbt{Khai thác đồ thị sóng và quan hệ pha giữa các điểm}
\begin{ex}
% ID: L11C2B8-16-TL
% Mức: VD
Một học sinh giải bài sau nhưng bỏ qua điều kiện vật lí. Hãy sửa cách làm và trình bày lời giải đúng: Một ảnh sóng có bước sóng $80\,\text{cm}$. Xác định ba vị trí gần nhất tính từ $O$ theo chiều truyền mà phần tử ngược pha với O.
\choiceTF

\loigiai{
Cách giải đúng: Điểm ngược pha thỏa $d=(k+1/2)\lambda$. Ba vị trí gần nhất là $40\,\text{cm}$, $120\,\text{cm}$ và $200\,\text{cm}$. Sai lầm cần tránh là dùng hệ thức khi chưa xác định điều kiện biên, môi trường hoặc đơn vị.
}
\end{ex}

% ===== Câu 63 =====
% ID: L11C2B8-16-TLN
% Mức: VD
\dangbt{Nhận biết bản chất và sự truyền sóng cơ}
\begin{ex}
% ID: L11C2B8-16-TLN
% Mức: VD
Từ kết quả của tình huống sau, nếu đại lượng cần tìm được nhân với hệ số 2 thì giá trị mới bằng bao nhiêu? Một sóng truyền quãng đường $18\,\text{m}$ trong $6\,\text{s}$. Tốc độ truyền sóng tính theo m/s bằng bao nhiêu?
\shortans{6}
\loigiai{
Xác định đúng dữ kiện và đơn vị, áp dụng hệ thức của dạng bài rồi tính được kết quả 6.
}
\end{ex}

% ===== Câu 64 =====
% ID: L11C2B8-16-TN
% Mức: TH
\dangbt{Khai thác đồ thị sóng và quan hệ pha giữa các điểm}
\begin{ex}
% ID: L11C2B8-16-TN
% Mức: TH
Tại một thời điểm, khoảng cách theo phương truyền giữa một đỉnh sóng và đáy gần nhất là
\choice
{\True $\lambda/2$}
{$\lambda$}
{$\lambda/4$}
{$2\lambda$}
\loigiai{
Đáp án A. $\lambda/2$. Đối chiếu định nghĩa, hệ thức hoặc dữ kiện của bài toán cho kết luận trên.
}
\end{ex}

% ===== Câu 65 =====
% ID: L11C2B8-17-DS
% Mức: TH
\dangbt{Phân tích trạng thái và chiều truyền sóng từ dữ kiện quan sát}
\begin{ex}
% ID: L11C2B8-17-DS
% Mức: TH
Trong chuyên đề “Phân tích trạng thái và chiều truyền sóng từ dữ kiện quan sát”, Ảnh chụp sóng cho thấy hai đỉnh liên tiếp cách nhau $40\,\text{cm}$.
\choiceTF
{\True Bước sóng bằng $40\,\text{cm}$.}
{\True Đỉnh và đáy gần nhất cách nhau $20\,\text{cm}$.}
{\True Hai điểm cách nhau $10\,\text{cm}$ vuông pha.}
{Hai điểm cách nhau $30\,\text{cm}$ cùng pha.}
\loigiai{
\begin{itemchoice}
\item (Đúng) Bước sóng bằng $40\,\text{cm}$. (Vì): Bước sóng bằng $40\,\text{cm}$. b) Đ: Đỉnh và đáy gần nhất cách nhau $20\,\text{cm}$. c) Đ: Hai điểm cách nhau $10\,\text{cm}$ vuông pha. d) S: Hai điểm cách nhau $30\,\text{cm}$ cùng pha
\item (Đúng) Đỉnh và đáy gần nhất cách nhau $20\,\text{cm}$.
\item (Đúng) Hai điểm cách nhau $10\,\text{cm}$ vuông pha.
\item (Sai) Hai điểm cách nhau $30\,\text{cm}$ cùng pha.
\end{itemchoice}
}
\end{ex}

% ===== Câu 66 =====
% ID: L11C2B8-17-TL
% Mức: VD
\dangbt{Khai thác đồ thị sóng và quan hệ pha giữa các điểm}
\begin{ex}
% ID: L11C2B8-17-TL
% Mức: VD
Mở rộng tình huống sau thành một ví dụ thực tế và trình bày phương pháp giải: Giải thích vì sao đường cong hình sin tại một thời điểm không phải là quỹ đạo chuyển động của một phần tử môi trường.
\choiceTF

\loigiai{
Đường cong biểu diễn li độ của nhiều phần tử ở các vị trí khác nhau tại cùng thời điểm; mỗi phần tử chỉ dao động quanh vị trí cân bằng riêng. Có thể kiểm chứng bằng cách thay kết quả vào dữ kiện ban đầu và đối chiếu với hiện tượng thực tế.
}
\end{ex}

% ===== Câu 67 =====
% ID: L11C2B8-17-TLN
% Mức: VD
\dangbt{Nhận biết bản chất và sự truyền sóng cơ}
\begin{ex}
% ID: L11C2B8-17-TLN
% Mức: VD
Từ kết quả của tình huống sau, nếu đại lượng cần tìm được nhân với hệ số 0.5 thì giá trị mới bằng bao nhiêu? Sóng truyền $24\,\text{m}$ trong $8\,\text{s}$. Tốc độ truyền sóng tính theo m/s bằng bao nhiêu?
\shortans{01/05/2026}
\loigiai{
Xác định đúng dữ kiện và đơn vị, áp dụng hệ thức của dạng bài rồi tính được kết quả 1.5.
}
\end{ex}

% ===== Câu 68 =====
% ID: L11C2B8-17-TN
% Mức: TH
\dangbt{Khai thác đồ thị sóng và quan hệ pha giữa các điểm}
\begin{ex}
% ID: L11C2B8-17-TN
% Mức: TH
Hai điểm cách nhau $0,6\,\text{m}$ có độ lệch pha $\pi$ rad và là hai điểm ngược pha gần nhau nhất. Bước sóng tính theo m bằng bao nhiêu? Chọn giá trị đúng.
\choice
{02/02/2026}
{\True 01/02/2026}
{02/04/2026}
{0.6}
\loigiai{
Đáp án B. 1.2. Đối chiếu định nghĩa, hệ thức hoặc dữ kiện của bài toán cho kết luận trên.
}
\end{ex}

% ===== Câu 69 =====
% ID: L11C2B8-18-DS
% Mức: TH
\dangbt{Phân tích trạng thái và chiều truyền sóng từ dữ kiện quan sát}
\begin{ex}
% ID: L11C2B8-18-DS
% Mức: TH
Trong chuyên đề “Phân tích trạng thái và chiều truyền sóng từ dữ kiện quan sát”, Sóng truyền theo chiều dương Ox.
\choiceTF
{\True Điểm gần nguồn hơn dao động sớm pha hơn điểm xa nguồn.}
{\True Pha giảm khi tọa độ tăng tại cùng thời điểm.}
{\True Hai điểm cách nhau một bước sóng có cùng trạng thái dao động.}
{Mọi điểm trên phương truyền đều có cùng pha.}
\loigiai{
\begin{itemchoice}
\item (Đúng) Điểm gần nguồn hơn dao động sớm pha hơn điểm xa nguồn. (Vì): iểm gần nguồn hơn dao động sớm pha hơn điểm xa nguồn. b) Đ: Pha giảm khi tọa độ tăng tại cùng thời điểm. c) Đ: Hai điểm cách nhau một bước sóng có cùng trạng thái dao động. d) S: Mọi điểm trên phương truyền đều có cùng pha
\item (Đúng) Pha giảm khi tọa độ tăng tại cùng thời điểm.
\item (Đúng) Hai điểm cách nhau một bước sóng có cùng trạng thái dao động.
\item (Sai) Mọi điểm trên phương truyền đều có cùng pha.
\end{itemchoice}
}
\end{ex}

% ===== Câu 70 =====
% ID: L11C2B8-18-TL
% Mức: TH
\dangbt{Nhận biết bản chất và sự truyền sóng cơ}
\begin{ex}
% ID: L11C2B8-18-TL
% Mức: TH
Giải thích vì sao sóng cơ truyền năng lượng nhưng không mang các phần tử môi trường đi theo.
\choiceTF

\loigiai{
Các phần tử môi trường chỉ dao động quanh vị trí cân bằng; sự liên kết giữa chúng truyền trạng thái dao động và năng lượng từ phần tử này sang phần tử khác.
}
\end{ex}

% ===== Câu 71 =====
% ID: L11C2B8-18-TLN
% Mức: VD
\begin{ex}
% ID: L11C2B8-18-TLN
% Mức: VD
Từ kết quả của tình huống sau, nếu đại lượng cần tìm được nhân với hệ số 1.5 thì giá trị mới bằng bao nhiêu? Một xung sóng đi được $15\,\text{m}$ với tốc độ $5\,\text{m}$ /s. Thời gian truyền tính theo s bằng bao nhiêu?
\shortans{04/05/2026}
\loigiai{
Xác định đúng dữ kiện và đơn vị, áp dụng hệ thức của dạng bài rồi tính được kết quả 4.5.
}
\end{ex}

% ===== Câu 72 =====
% ID: L11C2B8-18-TN
% Mức: TH
\dangbt{Khai thác đồ thị sóng và quan hệ pha giữa các điểm}
\begin{ex}
% ID: L11C2B8-18-TN
% Mức: TH
Bước sóng là $2,4\,\text{m}$. Khoảng cách gần nhất giữa hai điểm vuông pha tính theo m bằng bao nhiêu? Chọn giá trị đúng.
\choice
{0.3}
{01/06/2026}
{\True 0.6}
{01/02/2026}
\loigiai{
Đáp án C. 0.6. Đối chiếu định nghĩa, hệ thức hoặc dữ kiện của bài toán cho kết luận trên.
}
\end{ex}

% ===== Câu 73 =====
% ID: L11C2B8-19-DS
% Mức: VD
\dangbt{Phân tích trạng thái và chiều truyền sóng từ dữ kiện quan sát}
\begin{ex}
% ID: L11C2B8-19-DS
% Mức: VD
Trong chuyên đề “Phân tích trạng thái và chiều truyền sóng từ dữ kiện quan sát”, Xét các nhận định tổng hợp thuộc dạng “Phân tích trạng thái và chiều truyền sóng từ dữ kiện quan sát”.
\choiceTF
{\True một bước sóng.}
{\True ngược pha.}
{\True $\lambda/2$.}
{nửa bước sóng.}
\loigiai{
\begin{itemchoice}
\item (Đúng) một bước sóng. (Vì): một bước sóng. b) Đ: ngược pha. c) Đ: $\lambda/2$. d) S: nửa bước sóng
\item (Đúng) ngược pha.
\item (Đúng) $\lambda/2$.
\item (Sai) nửa bước sóng.
\end{itemchoice}
}
\end{ex}

% ===== Câu 74 =====
% ID: L11C2B8-19-TL
% Mức: TH
\dangbt{Nhận biết bản chất và sự truyền sóng cơ}
\begin{ex}
% ID: L11C2B8-19-TL
% Mức: TH
Nêu hai điều kiện cần để hình thành sóng cơ từ một nguồn dao động.
\choiceTF

\loigiai{
Cần có nguồn dao động và môi trường vật chất có lực liên kết giữa các phần tử để truyền dao động.
}
\end{ex}

% ===== Câu 75 =====
% ID: L11C2B8-19-TLN
% Mức: TH
\dangbt{Phân tích trạng thái và chiều truyền sóng từ dữ kiện quan sát}
\begin{ex}
% ID: L11C2B8-19-TLN
% Mức: TH
Trong chuyên đề “Phân tích trạng thái và chiều truyền sóng từ dữ kiện quan sát”, Hai điểm cách nhau $0,6\,\text{m}$ có độ lệch pha $\pi$ rad và là hai điểm ngược pha gần nhau nhất. Bước sóng tính theo m bằng bao nhiêu?
\shortans{01/02/2026}
\loigiai{
Xác định đúng dữ kiện và đơn vị, áp dụng hệ thức của dạng bài rồi tính được kết quả 1.2.
}
\end{ex}

% ===== Câu 76 =====
% ID: L11C2B8-19-TN
% Mức: TH
\dangbt{Khai thác đồ thị sóng và quan hệ pha giữa các điểm}
\begin{ex}
% ID: L11C2B8-19-TN
% Mức: TH
Hai đỉnh liên tiếp trên ảnh sóng cách nhau $75\,\text{cm}$. Khoảng cách từ một đỉnh đến đáy gần nhất tính theo cm bằng bao nhiêu? Chọn giá trị đúng.
\choice
{75}
{18.75}
{38.5}
{\True 37.5}
\loigiai{
Đáp án D. 37.5. Đối chiếu định nghĩa, hệ thức hoặc dữ kiện của bài toán cho kết luận trên.
}
\end{ex}

% ===== Câu 77 =====
% ID: L11C2B8-20-DS
% Mức: VD
\dangbt{Phân tích trạng thái và chiều truyền sóng từ dữ kiện quan sát}
\begin{ex}
% ID: L11C2B8-20-DS
% Mức: VD
Trong chuyên đề “Phân tích trạng thái và chiều truyền sóng từ dữ kiện quan sát”, Đánh giá cách xử lí các dữ kiện định lượng của dạng “Phân tích trạng thái và chiều truyền sóng từ dữ kiện quan sát”.
\choiceTF
{\True Kết quả của bài toán thứ nhất là 1.2.}
{\True Kết quả của bài toán thứ hai là 0.6.}
{\True Kết quả của bài toán thứ ba là 37.5.}
{Có thể bỏ qua hoàn toàn đơn vị và điều kiện vật lí khi tính toán.}
\loigiai{
\begin{itemchoice}
\item (Đúng) Kết quả của bài toán thứ nhất là 1.2. (Vì): Kết quả của bài toán thứ nhất là 1.2. b) Đ: Kết quả của bài toán thứ hai là 0.6. c) Đ: Kết quả của bài toán thứ ba là 37.5. d) S: Có thể bỏ qua hoàn toàn đơn vị và điều kiện vật lí khi tính toán
\item (Đúng) Kết quả của bài toán thứ hai là 0.6.
\item (Đúng) Kết quả của bài toán thứ ba là 37.5.
\item (Sai) Có thể bỏ qua hoàn toàn đơn vị và điều kiện vật lí khi tính toán.
\end{itemchoice}
}
\end{ex}

% ===== Câu 78 =====
% ID: L11C2B8-20-TL
% Mức: VD
\dangbt{Nhận biết bản chất và sự truyền sóng cơ}
\begin{ex}
% ID: L11C2B8-20-TL
% Mức: VD
Giải thích vì sao không thể nghe trực tiếp tiếng chuông đặt trong một bình đã hút hết không khí.
\choiceTF

\loigiai{
Âm là sóng cơ nên cần môi trường vật chất để truyền; trong chân không không có phần tử vật chất truyền dao động đến tai.
}
\end{ex}

% ===== Câu 79 =====
% ID: L11C2B8-20-TLN
% Mức: TH
\dangbt{Phân tích trạng thái và chiều truyền sóng từ dữ kiện quan sát}
\begin{ex}
% ID: L11C2B8-20-TLN
% Mức: TH
Trong chuyên đề “Phân tích trạng thái và chiều truyền sóng từ dữ kiện quan sát”, Bước sóng là $2,4\,\text{m}$. Khoảng cách gần nhất giữa hai điểm vuông pha tính theo m bằng bao nhiêu?
\shortans{0.6}
\loigiai{
Xác định đúng dữ kiện và đơn vị, áp dụng hệ thức của dạng bài rồi tính được kết quả 0.6.
}
\end{ex}

% ===== Câu 80 =====
% ID: L11C2B8-20-TN
% Mức: VD
\dangbt{Khai thác đồ thị sóng và quan hệ pha giữa các điểm}
\begin{ex}
% ID: L11C2B8-20-TN
% Mức: VD
Kết luận nào phù hợp nhất khi giải bài toán sau: Chứng minh hai điểm cách nhau một số nguyên lần bước sóng dao động cùng pha.
\choice
{\True Vì $\Delta\varphi=2\pi d/\lambda$. Với $d=k\lambda$ thì $\Delta\varphi=2k\pi$, nên hai điểm cùng pha.}
{Không cần sử dụng định nghĩa hay dữ kiện đã cho.}
{Đại lượng cần tìm luôn bằng 0 trong mọi trường hợp.}
{Kết quả không phụ thuộc môi trường, tần số hoặc điều kiện biên.}
\loigiai{
Đáp án A. Vì $\Delta\varphi=2\pi d/\lambda$. Với $d=k\lambda$ thì $\Delta\varphi=2k\pi$, nên hai điểm cùng pha.. Đối chiếu định nghĩa, hệ thức hoặc dữ kiện của bài toán cho kết luận trên.
}
\end{ex}

% ===== Câu 81 =====
% ID: L11C2B8-21-DS
% Mức: TH
\dangbt{Vận dụng khái niệm sóng cơ trong các hiện tượng thực tế}
\begin{ex}
% ID: L11C2B8-21-DS
% Mức: TH
Trong chuyên đề “Vận dụng khái niệm sóng cơ trong các hiện tượng thực tế”, Xét một sóng cơ truyền trên mặt nước.
\choiceTF
{\True Nguồn sóng cung cấp dao động cho môi trường.}
{\True Các phần tử nước dao động quanh vị trí cân bằng.}
{\True Sóng mang năng lượng ra xa nguồn.}
{Các phần tử nước chuyển động cùng sóng đến bờ.}
\loigiai{
\begin{itemchoice}
\item (Đúng) Nguồn sóng cung cấp dao động cho môi trường. (Vì): Nguồn sóng cung cấp dao động cho môi trường. b) Đ: Các phần tử nước dao động quanh vị trí cân bằng. c) Đ: Sóng mang năng lượng ra xa nguồn. d) S: Các phần tử nước chuyển động cùng sóng đến bờ
\item (Đúng) Các phần tử nước dao động quanh vị trí cân bằng.
\item (Đúng) Sóng mang năng lượng ra xa nguồn.
\item (Sai) Các phần tử nước chuyển động cùng sóng đến bờ.
\end{itemchoice}
}
\end{ex}

% ===== Câu 82 =====
% ID: L11C2B8-21-TL
% Mức: VD
\dangbt{Nhận biết bản chất và sự truyền sóng cơ}
\begin{ex}
% ID: L11C2B8-21-TL
% Mức: VD
Phân tích các bước lập luận và chỉ ra điều kiện áp dụng trong bài toán: Giải thích vì sao sóng cơ truyền năng lượng nhưng không mang các phần tử môi trường đi theo.
\choiceTF

\loigiai{
Các phần tử môi trường chỉ dao động quanh vị trí cân bằng; sự liên kết giữa chúng truyền trạng thái dao động và năng lượng từ phần tử này sang phần tử khác. Cần nêu rõ điều kiện áp dụng, đổi đơn vị thống nhất và kiểm tra ý nghĩa vật lí của kết quả.
}
\end{ex}

% ===== Câu 83 =====
% ID: L11C2B8-21-TLN
% Mức: TH
\dangbt{Phân tích trạng thái và chiều truyền sóng từ dữ kiện quan sát}
\begin{ex}
% ID: L11C2B8-21-TLN
% Mức: TH
Trong chuyên đề “Phân tích trạng thái và chiều truyền sóng từ dữ kiện quan sát”, Hai đỉnh liên tiếp trên ảnh sóng cách nhau $75\,\text{cm}$. Khoảng cách từ một đỉnh đến đáy gần nhất tính theo cm bằng bao nhiêu?
\shortans{37.5}
\loigiai{
Xác định đúng dữ kiện và đơn vị, áp dụng hệ thức của dạng bài rồi tính được kết quả 37.5.
}
\end{ex}

% ===== Câu 84 =====
% ID: L11C2B8-21-TN
% Mức: VD
\dangbt{Khai thác đồ thị sóng và quan hệ pha giữa các điểm}
\begin{ex}
% ID: L11C2B8-21-TN
% Mức: VD
Kết luận nào phù hợp nhất khi giải bài toán sau: Một ảnh sóng có bước sóng $80\,\text{cm}$. Xác định ba vị trí gần nhất tính từ $O$ theo chiều truyền mà phần tử ngược pha với O.
\choice
{Kết quả không phụ thuộc môi trường, tần số hoặc điều kiện biên.}
{\True Điểm ngược pha thỏa $d=(k+1/2)\lambda$. Ba vị trí gần nhất là $40\,\text{cm}$, $120\,\text{cm}$ và $200\,\text{cm}$.}
{Không cần sử dụng định nghĩa hay dữ kiện đã cho.}
{Đại lượng cần tìm luôn bằng 0 trong mọi trường hợp.}
\loigiai{
Đáp án B. Điểm ngược pha thỏa $d=(k+1/2)\lambda$. Ba vị trí gần nhất là $40\,\text{cm}$, $120\,\text{cm}$ và $200\,\text{cm}$.. Đối chiếu định nghĩa, hệ thức hoặc dữ kiện của bài toán cho kết luận trên.
}
\end{ex}

% ===== Câu 85 =====
% ID: L11C2B8-22-DS
% Mức: TH
\dangbt{Vận dụng khái niệm sóng cơ trong các hiện tượng thực tế}
\begin{ex}
% ID: L11C2B8-22-DS
% Mức: TH
Trong chuyên đề “Vận dụng khái niệm sóng cơ trong các hiện tượng thực tế”, Một loa phát âm trong phòng kín.
\choiceTF
{\True Âm là sóng cơ.}
{\True Không khí là môi trường truyền âm.}
{Âm truyền được trong chân không.}
{\True Khi âm tới tai, màng nhĩ nhận năng lượng.}
\loigiai{
\begin{itemchoice}
\item (Đúng) Âm là sóng cơ. (Vì): Âm là sóng cơ. b) Đ: Không khí là môi trường truyền âm. c) S: Âm truyền được trong chân không. d) Đ: Khi âm tới tai, màng nhĩ nhận năng lượng
\item (Đúng) Không khí là môi trường truyền âm.
\item (Sai) Âm truyền được trong chân không.
\item (Đúng) Khi âm tới tai, màng nhĩ nhận năng lượng.
\end{itemchoice}
}
\end{ex}

% ===== Câu 86 =====
% ID: L11C2B8-22-TL
% Mức: VD
\dangbt{Nhận biết bản chất và sự truyền sóng cơ}
\begin{ex}
% ID: L11C2B8-22-TL
% Mức: VD
Một học sinh giải bài sau nhưng bỏ qua điều kiện vật lí. Hãy sửa cách làm và trình bày lời giải đúng: Nêu hai điều kiện cần để hình thành sóng cơ từ một nguồn dao động.
\choiceTF

\loigiai{
Cách giải đúng: Cần có nguồn dao động và môi trường vật chất có lực liên kết giữa các phần tử để truyền dao động. Sai lầm cần tránh là dùng hệ thức khi chưa xác định điều kiện biên, môi trường hoặc đơn vị.
}
\end{ex}

% ===== Câu 87 =====
% ID: L11C2B8-22-TLN
% Mức: VD
\dangbt{Phân tích trạng thái và chiều truyền sóng từ dữ kiện quan sát}
\begin{ex}
% ID: L11C2B8-22-TLN
% Mức: VD
Trong chuyên đề “Phân tích trạng thái và chiều truyền sóng từ dữ kiện quan sát”, Từ kết quả của tình huống sau, nếu đại lượng cần tìm được nhân với hệ số 2 thì giá trị mới bằng bao nhiêu? Hai điểm cách nhau $0,6\,\text{m}$ có độ lệch pha $\pi$ rad và là hai điểm ngược pha gần nhau nhất. Bước sóng tính theo m bằng bao nhiêu?
\shortans{02/04/2026}
\loigiai{
Xác định đúng dữ kiện và đơn vị, áp dụng hệ thức của dạng bài rồi tính được kết quả 2.4.
}
\end{ex}

% ===== Câu 88 =====
% ID: L11C2B8-22-TN
% Mức: VD
\dangbt{Khai thác đồ thị sóng và quan hệ pha giữa các điểm}
\begin{ex}
% ID: L11C2B8-22-TN
% Mức: VD
Kết luận nào phù hợp nhất khi giải bài toán sau: Giải thích vì sao đường cong hình sin tại một thời điểm không phải là quỹ đạo chuyển động của một phần tử môi trường.
\choice
{Đại lượng cần tìm luôn bằng 0 trong mọi trường hợp.}
{Kết quả không phụ thuộc môi trường, tần số hoặc điều kiện biên.}
{\True Đường cong biểu diễn li độ của nhiều phần tử ở các vị trí khác nhau tại cùng thời điểm}
{Không cần sử dụng định nghĩa hay dữ kiện đã cho.}
\loigiai{
Đáp án C. Đường cong biểu diễn li độ của nhiều phần tử ở các vị trí khác nhau tại cùng thời điểm. Đối chiếu định nghĩa, hệ thức hoặc dữ kiện của bài toán cho kết luận trên.
}
\end{ex}

% ===== Câu 89 =====
% ID: L11C2B8-23-DS
% Mức: TH
\dangbt{Vận dụng khái niệm sóng cơ trong các hiện tượng thực tế}
\begin{ex}
% ID: L11C2B8-23-DS
% Mức: TH
Trong chuyên đề “Vận dụng khái niệm sóng cơ trong các hiện tượng thực tế”, Quan sát một xung truyền trên dây căng.
\choiceTF
{\True Biến dạng lan dọc theo dây.}
{\True Mỗi phần tử dây chỉ dao động quanh vị trí cân bằng.}
{Tốc độ truyền xung luôn bằng tốc độ dao động của phần tử dây.}
{\True Lực liên kết giữa các phần tử giúp truyền dao động.}
\loigiai{
\begin{itemchoice}
\item (Đúng) Biến dạng lan dọc theo dây. (Vì): Biến dạng lan dọc theo dây. b) Đ: Mỗi phần tử dây chỉ dao động quanh vị trí cân bằng. c) S: Tốc độ truyền xung luôn bằng tốc độ dao động của phần tử dây. d) Đ: Lực liên kết giữa các phần tử giúp truyền dao động
\item (Đúng) Mỗi phần tử dây chỉ dao động quanh vị trí cân bằng.
\item (Sai) Tốc độ truyền xung luôn bằng tốc độ dao động của phần tử dây.
\item (Đúng) Lực liên kết giữa các phần tử giúp truyền dao động.
\end{itemchoice}
}
\end{ex}

% ===== Câu 90 =====
% ID: L11C2B8-23-TL
% Mức: VD
\dangbt{Nhận biết bản chất và sự truyền sóng cơ}
\begin{ex}
% ID: L11C2B8-23-TL
% Mức: VD
Mở rộng tình huống sau thành một ví dụ thực tế và trình bày phương pháp giải: Giải thích vì sao không thể nghe trực tiếp tiếng chuông đặt trong một bình đã hút hết không khí.
\choiceTF

\loigiai{
Âm là sóng cơ nên cần môi trường vật chất để truyền; trong chân không không có phần tử vật chất truyền dao động đến tai. Có thể kiểm chứng bằng cách thay kết quả vào dữ kiện ban đầu và đối chiếu với hiện tượng thực tế.
}
\end{ex}

% ===== Câu 91 =====
% ID: L11C2B8-23-TLN
% Mức: VD
\dangbt{Phân tích trạng thái và chiều truyền sóng từ dữ kiện quan sát}
\begin{ex}
% ID: L11C2B8-23-TLN
% Mức: VD
Trong chuyên đề “Phân tích trạng thái và chiều truyền sóng từ dữ kiện quan sát”, Từ kết quả của tình huống sau, nếu đại lượng cần tìm được nhân với hệ số 0.5 thì giá trị mới bằng bao nhiêu? Bước sóng là $2,4\,\text{m}$. Khoảng cách gần nhất giữa hai điểm vuông pha tính theo m bằng bao nhiêu?
\shortans{0.3}
\loigiai{
Xác định đúng dữ kiện và đơn vị, áp dụng hệ thức của dạng bài rồi tính được kết quả 0.3.
}
\end{ex}

% ===== Câu 92 =====
% ID: L11C2B8-23-TN
% Mức: VD
\dangbt{Khai thác đồ thị sóng và quan hệ pha giữa các điểm}
\begin{ex}
% ID: L11C2B8-23-TN
% Mức: VD
Phương pháp phù hợp nhất cho dạng “Khai thác đồ thị sóng và quan hệ pha giữa các điểm” là
\choice
{chỉ thay số mà không cần xét điều kiện áp dụng}
{bỏ qua đơn vị và chiều truyền sóng}
{luôn coi mọi điểm dao động cùng pha}
{\True xác định dữ kiện, chọn đúng hệ thức hoặc định nghĩa, đổi đơn vị rồi kiểm tra kết quả}
\loigiai{
Đáp án D. xác định dữ kiện, chọn đúng hệ thức hoặc định nghĩa, đổi đơn vị rồi kiểm tra kết quả. Đối chiếu định nghĩa, hệ thức hoặc dữ kiện của bài toán cho kết luận trên.
}
\end{ex}

% ===== Câu 93 =====
% ID: L11C2B8-24-DS
% Mức: VD
\dangbt{Vận dụng khái niệm sóng cơ trong các hiện tượng thực tế}
\begin{ex}
% ID: L11C2B8-24-DS
% Mức: VD
Trong chuyên đề “Vận dụng khái niệm sóng cơ trong các hiện tượng thực tế”, Xét các nhận định tổng hợp thuộc dạng “Vận dụng khái niệm sóng cơ trong các hiện tượng thực tế”.
\choiceTF
{\True năng lượng và trạng thái dao động.}
{\True chân không.}
{\True dao động quanh vị trí cân bằng.}
{các phần tử vật chất.}
\loigiai{
\begin{itemchoice}
\item (Đúng) năng lượng và trạng thái dao động. (Vì): năng lượng và trạng thái dao động. b) Đ: chân không. c) Đ: dao động quanh vị trí cân bằng. d) S: các phần tử vật chất
\item (Đúng) chân không.
\item (Đúng) dao động quanh vị trí cân bằng.
\item (Sai) các phần tử vật chất.
\end{itemchoice}
}
\end{ex}

% ===== Câu 94 =====
% ID: L11C2B8-24-TL
% Mức: TH
\dangbt{Phân tích trạng thái và chiều truyền sóng từ dữ kiện quan sát}
\begin{ex}
% ID: L11C2B8-24-TL
% Mức: TH
Trong chuyên đề “Phân tích trạng thái và chiều truyền sóng từ dữ kiện quan sát”, Chứng minh hai điểm cách nhau một số nguyên lần bước sóng dao động cùng pha.
\choiceTF

\loigiai{
Vì $\Delta\varphi=2\pi d/\lambda$. Với $d=k\lambda$ thì $\Delta\varphi=2k\pi$, nên hai điểm cùng pha.
}
\end{ex}

% ===== Câu 95 =====
% ID: L11C2B8-24-TLN
% Mức: VD
\begin{ex}
% ID: L11C2B8-24-TLN
% Mức: VD
Trong chuyên đề “Phân tích trạng thái và chiều truyền sóng từ dữ kiện quan sát”, Từ kết quả của tình huống sau, nếu đại lượng cần tìm được nhân với hệ số 1.5 thì giá trị mới bằng bao nhiêu? Hai đỉnh liên tiếp trên ảnh sóng cách nhau $75\,\text{cm}$. Khoảng cách từ một đỉnh đến đáy gần nhất tính theo cm bằng bao nhiêu?
\shortans{56.25}
\loigiai{
Xác định đúng dữ kiện và đơn vị, áp dụng hệ thức của dạng bài rồi tính được kết quả 56.25.
}
\end{ex}

% ===== Câu 96 =====
% ID: L11C2B8-24-TN
% Mức: NB
\dangbt{Nhận biết bản chất và sự truyền sóng cơ}
\begin{ex}
% ID: L11C2B8-24-TN
% Mức: NB
Khi sóng cơ truyền qua một môi trường, đại lượng được truyền đi từ nơi này đến nơi khác là
\choice
{\True năng lượng và trạng thái dao động}
{các phần tử vật chất}
{khối lượng môi trường}
{điện tích của phần tử}
\loigiai{
Đáp án A. năng lượng và trạng thái dao động. Đối chiếu định nghĩa, hệ thức hoặc dữ kiện của bài toán cho kết luận trên.
}
\end{ex}

% ===== Câu 97 =====
% ID: L11C2B8-25-DS
% Mức: VD
\dangbt{Vận dụng khái niệm sóng cơ trong các hiện tượng thực tế}
\begin{ex}
% ID: L11C2B8-25-DS
% Mức: VD
Trong chuyên đề “Vận dụng khái niệm sóng cơ trong các hiện tượng thực tế”, Đánh giá cách xử lí các dữ kiện định lượng của dạng “Vận dụng khái niệm sóng cơ trong các hiện tượng thực tế”.
\choiceTF
{\True Kết quả của bài toán thứ nhất là 3.}
{\True Kết quả của bài toán thứ hai là 3.}
{\True Kết quả của bài toán thứ ba là 3.}
{Có thể bỏ qua hoàn toàn đơn vị và điều kiện vật lí khi tính toán.}
\loigiai{
\begin{itemchoice}
\item (Đúng) Kết quả của bài toán thứ nhất là 3. (Vì): Kết quả của bài toán thứ nhất là 3. b) Đ: Kết quả của bài toán thứ hai là 3. c) Đ: Kết quả của bài toán thứ ba là 3. d) S: Có thể bỏ qua hoàn toàn đơn vị và điều kiện vật lí khi tính toán
\item (Đúng) Kết quả của bài toán thứ hai là 3.
\item (Đúng) Kết quả của bài toán thứ ba là 3.
\item (Sai) Có thể bỏ qua hoàn toàn đơn vị và điều kiện vật lí khi tính toán.
\end{itemchoice}
}
\end{ex}

% ===== Câu 98 =====
% ID: L11C2B8-25-TL
% Mức: TH
\dangbt{Phân tích trạng thái và chiều truyền sóng từ dữ kiện quan sát}
\begin{ex}
% ID: L11C2B8-25-TL
% Mức: TH
Trong chuyên đề “Phân tích trạng thái và chiều truyền sóng từ dữ kiện quan sát”, Một ảnh sóng có bước sóng $80\,\text{cm}$. Xác định ba vị trí gần nhất tính từ $O$ theo chiều truyền mà phần tử ngược pha với O.
\choiceTF

\loigiai{
Điểm ngược pha thỏa $d=(k+1/2)\lambda$. Ba vị trí gần nhất là $40\,\text{cm}$, $120\,\text{cm}$ và $200\,\text{cm}$.
}
\end{ex}

% ===== Câu 99 =====
% ID: L11C2B8-25-TLN
% Mức: TH
\dangbt{Vận dụng khái niệm sóng cơ trong các hiện tượng thực tế}
\begin{ex}
% ID: L11C2B8-25-TLN
% Mức: TH
Trong chuyên đề “Vận dụng khái niệm sóng cơ trong các hiện tượng thực tế”, Một sóng truyền quãng đường $18\,\text{m}$ trong $6\,\text{s}$. Tốc độ truyền sóng tính theo m/s bằng bao nhiêu?
\shortans{3}
\loigiai{
Xác định đúng dữ kiện và đơn vị, áp dụng hệ thức của dạng bài rồi tính được kết quả 3.
}
\end{ex}

% ===== Câu 100 =====
% ID: L11C2B8-25-TN
% Mức: NB
\dangbt{Nhận biết bản chất và sự truyền sóng cơ}
\begin{ex}
% ID: L11C2B8-25-TN
% Mức: NB
Sóng cơ không thể truyền trong
\choice
{không khí}
{\True chân không}
{thép}
{nước}
\loigiai{
Đáp án B. chân không. Đối chiếu định nghĩa, hệ thức hoặc dữ kiện của bài toán cho kết luận trên.
}
\end{ex}

% ===== Câu 101 =====
% ID: L11C2B8-26-DS
% Mức: TH
\dangbt{Xác định bước sóng, chu kì, tần số và tốc độ truyền sóng}
\begin{ex}
% ID: L11C2B8-26-DS
% Mức: TH
Một sóng có $f=10$ Hz và $v=30$ m/s.
\choiceTF
{\True Chu kì bằng $0,1\,\text{s}$.}
{\True Bước sóng bằng $3\,\text{m}$.}
{\True Trong $2\,\text{s}$ sóng truyền được $60\,\text{m}$.}
{Tăng tần số trong cùng môi trường làm tốc độ truyền sóng tăng tương ứng.}
\loigiai{
\begin{itemchoice}
\item (Đúng) Chu kì bằng $0,1\,\text{s}$. (Vì): Chu kì bằng $0,1\,\text{s}$. b) Đ: Bước sóng bằng $3\,\text{m}$. c) Đ: Trong $2\,\text{s}$ sóng truyền được $60\,\text{m}$. d) S: Tăng tần số trong cùng môi trường làm tốc độ truyền sóng tăng tương ứng
\item (Đúng) Bước sóng bằng $3\,\text{m}$.
\item (Đúng) Trong $2\,\text{s}$ sóng truyền được $60\,\text{m}$.
\item (Sai) Tăng tần số trong cùng môi trường làm tốc độ truyền sóng tăng tương ứng.
\end{itemchoice}
}
\end{ex}

% ===== Câu 102 =====
% ID: L11C2B8-26-TL
% Mức: VD
\dangbt{Phân tích trạng thái và chiều truyền sóng từ dữ kiện quan sát}
\begin{ex}
% ID: L11C2B8-26-TL
% Mức: VD
Trong chuyên đề “Phân tích trạng thái và chiều truyền sóng từ dữ kiện quan sát”, Giải thích vì sao đường cong hình sin tại một thời điểm không phải là quỹ đạo chuyển động của một phần tử môi trường.
\choiceTF

\loigiai{
Đường cong biểu diễn li độ của nhiều phần tử ở các vị trí khác nhau tại cùng thời điểm; mỗi phần tử chỉ dao động quanh vị trí cân bằng riêng.
}
\end{ex}

% ===== Câu 103 =====
% ID: L11C2B8-26-TLN
% Mức: TH
\dangbt{Vận dụng khái niệm sóng cơ trong các hiện tượng thực tế}
\begin{ex}
% ID: L11C2B8-26-TLN
% Mức: TH
Trong chuyên đề “Vận dụng khái niệm sóng cơ trong các hiện tượng thực tế”, Sóng truyền $24\,\text{m}$ trong $8\,\text{s}$. Tốc độ truyền sóng tính theo m/s bằng bao nhiêu?
\shortans{3}
\loigiai{
Xác định đúng dữ kiện và đơn vị, áp dụng hệ thức của dạng bài rồi tính được kết quả 3.
}
\end{ex}

% ===== Câu 104 =====
% ID: L11C2B8-26-TN
% Mức: TH
\dangbt{Nhận biết bản chất và sự truyền sóng cơ}
\begin{ex}
% ID: L11C2B8-26-TN
% Mức: TH
Một nút chai nổi trên mặt hồ khi có sóng truyền qua chủ yếu
\choice
{chuyển thẳng đến bờ}
{đứng yên tuyệt đối}
{\True dao động quanh vị trí cân bằng}
{trôi theo đúng tốc độ truyền sóng}
\loigiai{
Đáp án C. dao động quanh vị trí cân bằng. Đối chiếu định nghĩa, hệ thức hoặc dữ kiện của bài toán cho kết luận trên.
}
\end{ex}

% ===== Câu 105 =====
% ID: L11C2B8-27-DS
% Mức: TH
\dangbt{Xác định bước sóng, chu kì, tần số và tốc độ truyền sóng}
\begin{ex}
% ID: L11C2B8-27-DS
% Mức: TH
Trên ảnh sóng, 6 đỉnh liên tiếp trải trên đoạn $10\,\text{m}$; chu kì là $0,5\,\text{s}$.
\choiceTF
{\True Có 5 bước sóng trên đoạn đã cho.}
{\True Bước sóng bằng $2\,\text{m}$.}
{\True Tần số bằng $2\,\text{Hz}$.}
{Tốc độ truyền sóng bằng $1\,\text{m}$ /s.}
\loigiai{
\begin{itemchoice}
\item (Đúng) Có 5 bước sóng trên đoạn đã cho. (Vì): Có 5 bước sóng trên đoạn đã cho. b) Đ: Bước sóng bằng $2\,\text{m}$. c) Đ: Tần số bằng $2\,\text{Hz}$. d) S: Tốc độ truyền sóng bằng $1\,\text{m}$ /s
\item (Đúng) Bước sóng bằng $2\,\text{m}$.
\item (Đúng) Tần số bằng $2\,\text{Hz}$.
\item (Sai) Tốc độ truyền sóng bằng $1\,\text{m}$ /s.
\end{itemchoice}
}
\end{ex}

% ===== Câu 106 =====
% ID: L11C2B8-27-TL
% Mức: VD
\dangbt{Phân tích trạng thái và chiều truyền sóng từ dữ kiện quan sát}
\begin{ex}
% ID: L11C2B8-27-TL
% Mức: VD
Trong chuyên đề “Phân tích trạng thái và chiều truyền sóng từ dữ kiện quan sát”, Phân tích các bước lập luận và chỉ ra điều kiện áp dụng trong bài toán: Chứng minh hai điểm cách nhau một số nguyên lần bước sóng dao động cùng pha.
\choiceTF

\loigiai{
Vì $\Delta\varphi=2\pi d/\lambda$. Với $d=k\lambda$ thì $\Delta\varphi=2k\pi$, nên hai điểm cùng pha. Cần nêu rõ điều kiện áp dụng, đổi đơn vị thống nhất và kiểm tra ý nghĩa vật lí của kết quả.
}
\end{ex}

% ===== Câu 107 =====
% ID: L11C2B8-27-TLN
% Mức: TH
\dangbt{Vận dụng khái niệm sóng cơ trong các hiện tượng thực tế}
\begin{ex}
% ID: L11C2B8-27-TLN
% Mức: TH
Trong chuyên đề “Vận dụng khái niệm sóng cơ trong các hiện tượng thực tế”, Một xung sóng đi được $15\,\text{m}$ với tốc độ $5\,\text{m}$ /s. Thời gian truyền tính theo s bằng bao nhiêu?
\shortans{3}
\loigiai{
Xác định đúng dữ kiện và đơn vị, áp dụng hệ thức của dạng bài rồi tính được kết quả 3.
}
\end{ex}

% ===== Câu 108 =====
% ID: L11C2B8-27-TN
% Mức: TH
\dangbt{Nhận biết bản chất và sự truyền sóng cơ}
\begin{ex}
% ID: L11C2B8-27-TN
% Mức: TH
Một sóng truyền quãng đường $18\,\text{m}$ trong $6\,\text{s}$. Tốc độ truyền sóng tính theo m/s bằng bao nhiêu? Chọn giá trị đúng.
\choice
{6}
{01/05/2026}
{4}
{\True 3}
\loigiai{
Đáp án D. 3. Đối chiếu định nghĩa, hệ thức hoặc dữ kiện của bài toán cho kết luận trên.
}
\end{ex}

% ===== Câu 109 =====
% ID: L11C2B8-28-DS
% Mức: TH
\dangbt{Xác định bước sóng, chu kì, tần số và tốc độ truyền sóng}
\begin{ex}
% ID: L11C2B8-28-DS
% Mức: TH
Sóng truyền từ môi trường 1 sang môi trường 2, nguồn không đổi.
\choiceTF
{\True Tần số sóng không đổi.}
{\True Chu kì sóng không đổi.}
{\True Nếu tốc độ giảm thì bước sóng giảm.}
{Tích $v\lambda$ không đổi.}
\loigiai{
\begin{itemchoice}
\item (Đúng) Tần số sóng không đổi. (Vì): Tần số sóng không đổi. b) Đ: Chu kì sóng không đổi. c) Đ: Nếu tốc độ giảm thì bước sóng giảm. d) S: Tích $v\lambda$ không đổi
\item (Đúng) Chu kì sóng không đổi.
\item (Đúng) Nếu tốc độ giảm thì bước sóng giảm.
\item (Sai) Tích $v\lambda$ không đổi.
\end{itemchoice}
}
\end{ex}

% ===== Câu 110 =====
% ID: L11C2B8-28-TL
% Mức: VD
\dangbt{Phân tích trạng thái và chiều truyền sóng từ dữ kiện quan sát}
\begin{ex}
% ID: L11C2B8-28-TL
% Mức: VD
Trong chuyên đề “Phân tích trạng thái và chiều truyền sóng từ dữ kiện quan sát”, Một học sinh giải bài sau nhưng bỏ qua điều kiện vật lí. Hãy sửa cách làm và trình bày lời giải đúng: Một ảnh sóng có bước sóng $80\,\text{cm}$. Xác định ba vị trí gần nhất tính từ $O$ theo chiều truyền mà phần tử ngược pha với O.
\choiceTF

\loigiai{
Cách giải đúng: Điểm ngược pha thỏa $d=(k+1/2)\lambda$. Ba vị trí gần nhất là $40\,\text{cm}$, $120\,\text{cm}$ và $200\,\text{cm}$. Sai lầm cần tránh là dùng hệ thức khi chưa xác định điều kiện biên, môi trường hoặc đơn vị.
}
\end{ex}

% ===== Câu 111 =====
% ID: L11C2B8-28-TLN
% Mức: VD
\dangbt{Vận dụng khái niệm sóng cơ trong các hiện tượng thực tế}
\begin{ex}
% ID: L11C2B8-28-TLN
% Mức: VD
Trong chuyên đề “Vận dụng khái niệm sóng cơ trong các hiện tượng thực tế”, Từ kết quả của tình huống sau, nếu đại lượng cần tìm được nhân với hệ số 2 thì giá trị mới bằng bao nhiêu? Một sóng truyền quãng đường $18\,\text{m}$ trong $6\,\text{s}$. Tốc độ truyền sóng tính theo m/s bằng bao nhiêu?
\shortans{6}
\loigiai{
Xác định đúng dữ kiện và đơn vị, áp dụng hệ thức của dạng bài rồi tính được kết quả 6.
}
\end{ex}

% ===== Câu 112 =====
% ID: L11C2B8-28-TN
% Mức: TH
\dangbt{Nhận biết bản chất và sự truyền sóng cơ}
\begin{ex}
% ID: L11C2B8-28-TN
% Mức: TH
Sóng truyền $24\,\text{m}$ trong $8\,\text{s}$. Tốc độ truyền sóng tính theo m/s bằng bao nhiêu? Chọn giá trị đúng.
\choice
{\True 3}
{6}
{01/05/2026}
{4}
\loigiai{
Đáp án A. 3. Đối chiếu định nghĩa, hệ thức hoặc dữ kiện của bài toán cho kết luận trên.
}
\end{ex}

% ===== Câu 113 =====
% ID: L11C2B8-29-DS
% Mức: VD
\dangbt{Xác định bước sóng, chu kì, tần số và tốc độ truyền sóng}
\begin{ex}
% ID: L11C2B8-29-DS
% Mức: VD
Xét các nhận định tổng hợp thuộc dạng “Xác định bước sóng, chu kì, tần số và tốc độ truyền sóng”.
\choiceTF
{\True $4\,\text{m}$.}
{\True $2\,\text{m}$.}
{\True $4\,\text{m}$ /s.}
{$100\,\text{m}$.}
\loigiai{
\begin{itemchoice}
\item (Đúng) $4\,\text{m}$. (Vì): $4\,\text{m}$. b) Đ: $2\,\text{m}$. c) Đ: $4\,\text{m}$ /s. d) S: $100\,\text{m}$
\item (Đúng) $2\,\text{m}$.
\item (Đúng) $4\,\text{m}$ /s.
\item (Sai) $100\,\text{m}$.
\end{itemchoice}
}
\end{ex}

% ===== Câu 114 =====
% ID: L11C2B8-29-TL
% Mức: VD
\dangbt{Phân tích trạng thái và chiều truyền sóng từ dữ kiện quan sát}
\begin{ex}
% ID: L11C2B8-29-TL
% Mức: VD
Trong chuyên đề “Phân tích trạng thái và chiều truyền sóng từ dữ kiện quan sát”, Mở rộng tình huống sau thành một ví dụ thực tế và trình bày phương pháp giải: Giải thích vì sao đường cong hình sin tại một thời điểm không phải là quỹ đạo chuyển động của một phần tử môi trường.
\choiceTF

\loigiai{
Đường cong biểu diễn li độ của nhiều phần tử ở các vị trí khác nhau tại cùng thời điểm; mỗi phần tử chỉ dao động quanh vị trí cân bằng riêng. Có thể kiểm chứng bằng cách thay kết quả vào dữ kiện ban đầu và đối chiếu với hiện tượng thực tế.
}
\end{ex}

% ===== Câu 115 =====
% ID: L11C2B8-29-TLN
% Mức: VD
\dangbt{Vận dụng khái niệm sóng cơ trong các hiện tượng thực tế}
\begin{ex}
% ID: L11C2B8-29-TLN
% Mức: VD
Trong chuyên đề “Vận dụng khái niệm sóng cơ trong các hiện tượng thực tế”, Từ kết quả của tình huống sau, nếu đại lượng cần tìm được nhân với hệ số 0.5 thì giá trị mới bằng bao nhiêu? Sóng truyền $24\,\text{m}$ trong $8\,\text{s}$. Tốc độ truyền sóng tính theo m/s bằng bao nhiêu?
\shortans{01/05/2026}
\loigiai{
Xác định đúng dữ kiện và đơn vị, áp dụng hệ thức của dạng bài rồi tính được kết quả 1.5.
}
\end{ex}

% ===== Câu 116 =====
% ID: L11C2B8-29-TN
% Mức: TH
\dangbt{Nhận biết bản chất và sự truyền sóng cơ}
\begin{ex}
% ID: L11C2B8-29-TN
% Mức: TH
Một xung sóng đi được $15\,\text{m}$ với tốc độ $5\,\text{m}$ /s. Thời gian truyền tính theo s bằng bao nhiêu? Chọn giá trị đúng.
\choice
{4}
{\True 3}
{6}
{01/05/2026}
\loigiai{
Đáp án B. 3. Đối chiếu định nghĩa, hệ thức hoặc dữ kiện của bài toán cho kết luận trên.
}
\end{ex}

% ===== Câu 117 =====
% ID: L11C2B8-30-DS
% Mức: VD
\dangbt{Xác định bước sóng, chu kì, tần số và tốc độ truyền sóng}
\begin{ex}
% ID: L11C2B8-30-DS
% Mức: VD
Đánh giá cách xử lí các dữ kiện định lượng của dạng “Xác định bước sóng, chu kì, tần số và tốc độ truyền sóng”.
\choiceTF
{\True Kết quả của bài toán thứ nhất là 3.}
{\True Kết quả của bài toán thứ hai là 2.}
{\True Kết quả của bài toán thứ ba là 6.}
{Có thể bỏ qua hoàn toàn đơn vị và điều kiện vật lí khi tính toán.}
\loigiai{
\begin{itemchoice}
\item (Đúng) Kết quả của bài toán thứ nhất là 3. (Vì): Kết quả của bài toán thứ nhất là 3. b) Đ: Kết quả của bài toán thứ hai là 2. c) Đ: Kết quả của bài toán thứ ba là 6. d) S: Có thể bỏ qua hoàn toàn đơn vị và điều kiện vật lí khi tính toán
\item (Đúng) Kết quả của bài toán thứ hai là 2.
\item (Đúng) Kết quả của bài toán thứ ba là 6.
\item (Sai) Có thể bỏ qua hoàn toàn đơn vị và điều kiện vật lí khi tính toán.
\end{itemchoice}
}
\end{ex}

% ===== Câu 118 =====
% ID: L11C2B8-30-TL
% Mức: TH
\dangbt{Vận dụng khái niệm sóng cơ trong các hiện tượng thực tế}
\begin{ex}
% ID: L11C2B8-30-TL
% Mức: TH
Trong chuyên đề “Vận dụng khái niệm sóng cơ trong các hiện tượng thực tế”, Giải thích vì sao sóng cơ truyền năng lượng nhưng không mang các phần tử môi trường đi theo.
\choiceTF

\loigiai{
Các phần tử môi trường chỉ dao động quanh vị trí cân bằng; sự liên kết giữa chúng truyền trạng thái dao động và năng lượng từ phần tử này sang phần tử khác.
}
\end{ex}

% ===== Câu 119 =====
% ID: L11C2B8-30-TLN
% Mức: VD
\begin{ex}
% ID: L11C2B8-30-TLN
% Mức: VD
Trong chuyên đề “Vận dụng khái niệm sóng cơ trong các hiện tượng thực tế”, Từ kết quả của tình huống sau, nếu đại lượng cần tìm được nhân với hệ số 1.5 thì giá trị mới bằng bao nhiêu? Một xung sóng đi được $15\,\text{m}$ với tốc độ $5\,\text{m}$ /s. Thời gian truyền tính theo s bằng bao nhiêu?
\shortans{04/05/2026}
\loigiai{
Xác định đúng dữ kiện và đơn vị, áp dụng hệ thức của dạng bài rồi tính được kết quả 4.5.
}
\end{ex}

% ===== Câu 120 =====
% ID: L11C2B8-30-TN
% Mức: VD
\dangbt{Nhận biết bản chất và sự truyền sóng cơ}
\begin{ex}
% ID: L11C2B8-30-TN
% Mức: VD
Kết luận nào phù hợp nhất khi giải bài toán sau: Giải thích vì sao sóng cơ truyền năng lượng nhưng không mang các phần tử môi trường đi theo.
\choice
{Đại lượng cần tìm luôn bằng 0 trong mọi trường hợp.}
{Kết quả không phụ thuộc môi trường, tần số hoặc điều kiện biên.}
{\True Các phần tử môi trường chỉ dao động quanh vị trí cân bằng}
{Không cần sử dụng định nghĩa hay dữ kiện đã cho.}
\loigiai{
Đáp án C. Các phần tử môi trường chỉ dao động quanh vị trí cân bằng. Đối chiếu định nghĩa, hệ thức hoặc dữ kiện của bài toán cho kết luận trên.
}
\end{ex}

% ===== Câu 121 =====
% ID: L11C2B8-31-TL
% Mức: TH
\dangbt{Vận dụng khái niệm sóng cơ trong các hiện tượng thực tế}
\begin{ex}
% ID: L11C2B8-31-TL
% Mức: TH
Trong chuyên đề “Vận dụng khái niệm sóng cơ trong các hiện tượng thực tế”, Nêu hai điều kiện cần để hình thành sóng cơ từ một nguồn dao động.
\choiceTF

\loigiai{
Cần có nguồn dao động và môi trường vật chất có lực liên kết giữa các phần tử để truyền dao động.
}
\end{ex}

% ===== Câu 122 =====
% ID: L11C2B8-31-TLN
% Mức: TH
\dangbt{Xác định bước sóng, chu kì, tần số và tốc độ truyền sóng}
\begin{ex}
% ID: L11C2B8-31-TLN
% Mức: TH
Sóng có tốc độ $12\,\text{m}$ /s và tần số $4\,\text{Hz}$. Bước sóng tính theo m bằng bao nhiêu?
\shortans{3}
\loigiai{
Xác định đúng dữ kiện và đơn vị, áp dụng hệ thức của dạng bài rồi tính được kết quả 3.
}
\end{ex}

% ===== Câu 123 =====
% ID: L11C2B8-31-TN
% Mức: VD
\dangbt{Nhận biết bản chất và sự truyền sóng cơ}
\begin{ex}
% ID: L11C2B8-31-TN
% Mức: VD
Kết luận nào phù hợp nhất khi giải bài toán sau: Nêu hai điều kiện cần để hình thành sóng cơ từ một nguồn dao động.
\choice
{Không cần sử dụng định nghĩa hay dữ kiện đã cho.}
{Đại lượng cần tìm luôn bằng 0 trong mọi trường hợp.}
{Kết quả không phụ thuộc môi trường, tần số hoặc điều kiện biên.}
{\True Cần có nguồn dao động và môi trường vật chất có lực liên kết giữa các phần tử để truyền dao động.}
\loigiai{
Đáp án D. Cần có nguồn dao động và môi trường vật chất có lực liên kết giữa các phần tử để truyền dao động.. Đối chiếu định nghĩa, hệ thức hoặc dữ kiện của bài toán cho kết luận trên.
}
\end{ex}

% ===== Câu 124 =====
% ID: L11C2B8-32-TL
% Mức: VD
\dangbt{Vận dụng khái niệm sóng cơ trong các hiện tượng thực tế}
\begin{ex}
% ID: L11C2B8-32-TL
% Mức: VD
Trong chuyên đề “Vận dụng khái niệm sóng cơ trong các hiện tượng thực tế”, Giải thích vì sao không thể nghe trực tiếp tiếng chuông đặt trong một bình đã hút hết không khí.
\choiceTF

\loigiai{
Âm là sóng cơ nên cần môi trường vật chất để truyền; trong chân không không có phần tử vật chất truyền dao động đến tai.
}
\end{ex}

% ===== Câu 125 =====
% ID: L11C2B8-32-TLN
% Mức: TH
\dangbt{Xác định bước sóng, chu kì, tần số và tốc độ truyền sóng}
\begin{ex}
% ID: L11C2B8-32-TLN
% Mức: TH
Khoảng cách giữa 9 đỉnh sóng liên tiếp là $16\,\text{m}$. Bước sóng tính theo m bằng bao nhiêu?
\shortans{2}
\loigiai{
Xác định đúng dữ kiện và đơn vị, áp dụng hệ thức của dạng bài rồi tính được kết quả 2.
}
\end{ex}

% ===== Câu 126 =====
% ID: L11C2B8-32-TN
% Mức: VD
\dangbt{Nhận biết bản chất và sự truyền sóng cơ}
\begin{ex}
% ID: L11C2B8-32-TN
% Mức: VD
Kết luận nào phù hợp nhất khi giải bài toán sau: Giải thích vì sao không thể nghe trực tiếp tiếng chuông đặt trong một bình đã hút hết không khí.
\choice
{\True Âm là sóng cơ nên cần môi trường vật chất để truyền}
{Không cần sử dụng định nghĩa hay dữ kiện đã cho.}
{Đại lượng cần tìm luôn bằng 0 trong mọi trường hợp.}
{Kết quả không phụ thuộc môi trường, tần số hoặc điều kiện biên.}
\loigiai{
Đáp án A. Âm là sóng cơ nên cần môi trường vật chất để truyền. Đối chiếu định nghĩa, hệ thức hoặc dữ kiện của bài toán cho kết luận trên.
}
\end{ex}

% ===== Câu 127 =====
% ID: L11C2B8-33-TL
% Mức: VD
\dangbt{Vận dụng khái niệm sóng cơ trong các hiện tượng thực tế}
\begin{ex}
% ID: L11C2B8-33-TL
% Mức: VD
Trong chuyên đề “Vận dụng khái niệm sóng cơ trong các hiện tượng thực tế”, Phân tích các bước lập luận và chỉ ra điều kiện áp dụng trong bài toán: Giải thích vì sao sóng cơ truyền năng lượng nhưng không mang các phần tử môi trường đi theo.
\choiceTF

\loigiai{
Các phần tử môi trường chỉ dao động quanh vị trí cân bằng; sự liên kết giữa chúng truyền trạng thái dao động và năng lượng từ phần tử này sang phần tử khác. Cần nêu rõ điều kiện áp dụng, đổi đơn vị thống nhất và kiểm tra ý nghĩa vật lí của kết quả.
}
\end{ex}

% ===== Câu 128 =====
% ID: L11C2B8-33-TLN
% Mức: TH
\dangbt{Xác định bước sóng, chu kì, tần số và tốc độ truyền sóng}
\begin{ex}
% ID: L11C2B8-33-TLN
% Mức: TH
Sóng có bước sóng $1,5\,\text{m}$ và chu kì $0,25\,\text{s}$. Tốc độ truyền tính theo m/s bằng bao nhiêu?
\shortans{6}
\loigiai{
Xác định đúng dữ kiện và đơn vị, áp dụng hệ thức của dạng bài rồi tính được kết quả 6.
}
\end{ex}

% ===== Câu 129 =====
% ID: L11C2B8-33-TN
% Mức: VD
\dangbt{Nhận biết bản chất và sự truyền sóng cơ}
\begin{ex}
% ID: L11C2B8-33-TN
% Mức: VD
Phương pháp phù hợp nhất cho dạng “Nhận biết bản chất và sự truyền sóng cơ” là
\choice
{luôn coi mọi điểm dao động cùng pha}
{\True xác định dữ kiện, chọn đúng hệ thức hoặc định nghĩa, đổi đơn vị rồi kiểm tra kết quả}
{chỉ thay số mà không cần xét điều kiện áp dụng}
{bỏ qua đơn vị và chiều truyền sóng}
\loigiai{
Đáp án B. xác định dữ kiện, chọn đúng hệ thức hoặc định nghĩa, đổi đơn vị rồi kiểm tra kết quả. Đối chiếu định nghĩa, hệ thức hoặc dữ kiện của bài toán cho kết luận trên.
}
\end{ex}

% ===== Câu 130 =====
% ID: L11C2B8-34-TL
% Mức: VD
\dangbt{Vận dụng khái niệm sóng cơ trong các hiện tượng thực tế}
\begin{ex}
% ID: L11C2B8-34-TL
% Mức: VD
Trong chuyên đề “Vận dụng khái niệm sóng cơ trong các hiện tượng thực tế”, Một học sinh giải bài sau nhưng bỏ qua điều kiện vật lí. Hãy sửa cách làm và trình bày lời giải đúng: Nêu hai điều kiện cần để hình thành sóng cơ từ một nguồn dao động.
\choiceTF

\loigiai{
Cách giải đúng: Cần có nguồn dao động và môi trường vật chất có lực liên kết giữa các phần tử để truyền dao động. Sai lầm cần tránh là dùng hệ thức khi chưa xác định điều kiện biên, môi trường hoặc đơn vị.
}
\end{ex}

% ===== Câu 131 =====
% ID: L11C2B8-34-TLN
% Mức: VD
\dangbt{Xác định bước sóng, chu kì, tần số và tốc độ truyền sóng}
\begin{ex}
% ID: L11C2B8-34-TLN
% Mức: VD
Từ kết quả của tình huống sau, nếu đại lượng cần tìm được nhân với hệ số 2 thì giá trị mới bằng bao nhiêu? Sóng có tốc độ $12\,\text{m}$ /s và tần số $4\,\text{Hz}$. Bước sóng tính theo m bằng bao nhiêu?
\shortans{6}
\loigiai{
Xác định đúng dữ kiện và đơn vị, áp dụng hệ thức của dạng bài rồi tính được kết quả 6.
}
\end{ex}

% ===== Câu 132 =====
% ID: L11C2B8-34-TN
% Mức: NB
\dangbt{Phân tích trạng thái và chiều truyền sóng từ dữ kiện quan sát}
\begin{ex}
% ID: L11C2B8-34-TN
% Mức: NB
Trong chuyên đề “Phân tích trạng thái và chiều truyền sóng từ dữ kiện quan sát”, Hai điểm gần nhau nhất trên cùng phương truyền sóng dao động cùng pha cách nhau
\choice
{hai bước sóng}
{\True một bước sóng}
{nửa bước sóng}
{một phần tư bước sóng}
\loigiai{
Đáp án B. một bước sóng. Đối chiếu định nghĩa, hệ thức hoặc dữ kiện của bài toán cho kết luận trên.
}
\end{ex}

% ===== Câu 133 =====
% ID: L11C2B8-35-TL
% Mức: VD
\dangbt{Vận dụng khái niệm sóng cơ trong các hiện tượng thực tế}
\begin{ex}
% ID: L11C2B8-35-TL
% Mức: VD
Trong chuyên đề “Vận dụng khái niệm sóng cơ trong các hiện tượng thực tế”, Mở rộng tình huống sau thành một ví dụ thực tế và trình bày phương pháp giải: Giải thích vì sao không thể nghe trực tiếp tiếng chuông đặt trong một bình đã hút hết không khí.
\choiceTF

\loigiai{
Âm là sóng cơ nên cần môi trường vật chất để truyền; trong chân không không có phần tử vật chất truyền dao động đến tai. Có thể kiểm chứng bằng cách thay kết quả vào dữ kiện ban đầu và đối chiếu với hiện tượng thực tế.
}
\end{ex}

% ===== Câu 134 =====
% ID: L11C2B8-35-TLN
% Mức: VD
\dangbt{Xác định bước sóng, chu kì, tần số và tốc độ truyền sóng}
\begin{ex}
% ID: L11C2B8-35-TLN
% Mức: VD
Từ kết quả của tình huống sau, nếu đại lượng cần tìm được nhân với hệ số 0.5 thì giá trị mới bằng bao nhiêu? Khoảng cách giữa 9 đỉnh sóng liên tiếp là $16\,\text{m}$. Bước sóng tính theo m bằng bao nhiêu?
\shortans{1}
\loigiai{
Xác định đúng dữ kiện và đơn vị, áp dụng hệ thức của dạng bài rồi tính được kết quả 1.
}
\end{ex}

% ===== Câu 135 =====
% ID: L11C2B8-35-TN
% Mức: NB
\dangbt{Phân tích trạng thái và chiều truyền sóng từ dữ kiện quan sát}
\begin{ex}
% ID: L11C2B8-35-TN
% Mức: NB
Trong chuyên đề “Phân tích trạng thái và chiều truyền sóng từ dữ kiện quan sát”, Hai điểm cách nhau $\lambda/2$ trên cùng phương truyền sóng dao động
\choice
{vuông pha}
{không xác định}
{\True ngược pha}
{cùng pha}
\loigiai{
Đáp án C. ngược pha. Đối chiếu định nghĩa, hệ thức hoặc dữ kiện của bài toán cho kết luận trên.
}
\end{ex}

% ===== Câu 136 =====
% ID: L11C2B8-36-TL
% Mức: TH
\dangbt{Xác định bước sóng, chu kì, tần số và tốc độ truyền sóng}
\begin{ex}
% ID: L11C2B8-36-TL
% Mức: TH
Một sóng có tần số $8\,\text{Hz}$ truyền với tốc độ $24\,\text{m}$ /s. Tính chu kì, bước sóng và quãng đường sóng truyền trong $1,5\,\text{s}$.
\choiceTF

\loigiai{
$T=1/f=0,125$ s; $\lambda=v/f=3$ m; $s=vt=36$ m.
}
\end{ex}

% ===== Câu 137 =====
% ID: L11C2B8-36-TLN
% Mức: VD
\begin{ex}
% ID: L11C2B8-36-TLN
% Mức: VD
Từ kết quả của tình huống sau, nếu đại lượng cần tìm được nhân với hệ số 1.5 thì giá trị mới bằng bao nhiêu? Sóng có bước sóng $1,5\,\text{m}$ và chu kì $0,25\,\text{s}$. Tốc độ truyền tính theo m/s bằng bao nhiêu?
\shortans{9}
\loigiai{
Xác định đúng dữ kiện và đơn vị, áp dụng hệ thức của dạng bài rồi tính được kết quả 9.
}
\end{ex}

% ===== Câu 138 =====
% ID: L11C2B8-36-TN
% Mức: TH
\dangbt{Phân tích trạng thái và chiều truyền sóng từ dữ kiện quan sát}
\begin{ex}
% ID: L11C2B8-36-TN
% Mức: TH
Trong chuyên đề “Phân tích trạng thái và chiều truyền sóng từ dữ kiện quan sát”, Tại một thời điểm, khoảng cách theo phương truyền giữa một đỉnh sóng và đáy gần nhất là
\choice
{$\lambda$}
{$\lambda/4$}
{$2\lambda$}
{\True $\lambda/2$}
\loigiai{
Đáp án D. $\lambda/2$. Đối chiếu định nghĩa, hệ thức hoặc dữ kiện của bài toán cho kết luận trên.
}
\end{ex}

% ===== Câu 139 =====
% ID: L11C2B8-37-TL
% Mức: TH
\dangbt{Xác định bước sóng, chu kì, tần số và tốc độ truyền sóng}
\begin{ex}
% ID: L11C2B8-37-TL
% Mức: TH
Khoảng cách giữa đỉnh thứ nhất và đỉnh thứ mười một là $25\,\text{m}$. Sóng truyền với tốc độ $10\,\text{m}$ /s. Tính bước sóng và tần số.
\choiceTF

\loigiai{
Giữa 11 đỉnh có 10 bước sóng nên $\lambda=25/10=2,5$ m; $f=v/\lambda=4$ Hz.
}
\end{ex}

% ===== Câu 140 =====
% ID: L11C2B8-37-TN
% Mức: TH
\dangbt{Phân tích trạng thái và chiều truyền sóng từ dữ kiện quan sát}
\begin{ex}
% ID: L11C2B8-37-TN
% Mức: TH
Trong chuyên đề “Phân tích trạng thái và chiều truyền sóng từ dữ kiện quan sát”, Hai điểm cách nhau $0,6\,\text{m}$ có độ lệch pha $\pi$ rad và là hai điểm ngược pha gần nhau nhất. Bước sóng tính theo m bằng bao nhiêu? Chọn giá trị đúng.
\choice
{\True 01/02/2026}
{02/04/2026}
{0.6}
{02/02/2026}
\loigiai{
Đáp án A. 1.2. Đối chiếu định nghĩa, hệ thức hoặc dữ kiện của bài toán cho kết luận trên.
}
\end{ex}

% ===== Câu 141 =====
% ID: L11C2B8-38-TL
% Mức: VD
\dangbt{Xác định bước sóng, chu kì, tần số và tốc độ truyền sóng}
\begin{ex}
% ID: L11C2B8-38-TL
% Mức: VD
Một sóng đi từ môi trường $A$ sang $B$, tần số $20\,\text{Hz}$ không đổi. Biết tốc độ ở $A$ là $40\,\text{m}$ /s, ở $B$ là $30\,\text{m}$ /s. So sánh bước sóng ở hai môi trường.
\choiceTF

\loigiai{
$\lambda_A=2$ m, $\lambda_B=1,5$ m; bước sóng ở $B$ nhỏ hơn ở $A0,5\,\text{m}$.
}
\end{ex}

% ===== Câu 142 =====
% ID: L11C2B8-38-TN
% Mức: TH
\dangbt{Phân tích trạng thái và chiều truyền sóng từ dữ kiện quan sát}
\begin{ex}
% ID: L11C2B8-38-TN
% Mức: TH
Trong chuyên đề “Phân tích trạng thái và chiều truyền sóng từ dữ kiện quan sát”, Bước sóng là $2,4\,\text{m}$. Khoảng cách gần nhất giữa hai điểm vuông pha tính theo m bằng bao nhiêu? Chọn giá trị đúng.
\choice
{01/06/2026}
{\True 0.6}
{1.2}
{0.3}
\loigiai{
Đáp án B. 0.6. Đối chiếu định nghĩa, hệ thức hoặc dữ kiện của bài toán cho kết luận trên.
}
\end{ex}

% ===== Câu 143 =====
% ID: L11C2B8-39-TL
% Mức: VD
\dangbt{Xác định bước sóng, chu kì, tần số và tốc độ truyền sóng}
\begin{ex}
% ID: L11C2B8-39-TL
% Mức: VD
Phân tích các bước lập luận và chỉ ra điều kiện áp dụng trong bài toán: Một sóng có tần số $8\,\text{Hz}$ truyền với tốc độ $24\,\text{m}$ /s. Tính chu kì, bước sóng và quãng đường sóng truyền trong $1,5\,\text{s}$.
\choiceTF

\loigiai{
$T=1/f=0,125$ s; $\lambda=v/f=3$ m; $s=vt=36$ m. Cần nêu rõ điều kiện áp dụng, đổi đơn vị thống nhất và kiểm tra ý nghĩa vật lí của kết quả.
}
\end{ex}

% ===== Câu 144 =====
% ID: L11C2B8-39-TN
% Mức: TH
\dangbt{Phân tích trạng thái và chiều truyền sóng từ dữ kiện quan sát}
\begin{ex}
% ID: L11C2B8-39-TN
% Mức: TH
Trong chuyên đề “Phân tích trạng thái và chiều truyền sóng từ dữ kiện quan sát”, Hai đỉnh liên tiếp trên ảnh sóng cách nhau $75\,\text{cm}$. Khoảng cách từ một đỉnh đến đáy gần nhất tính theo cm bằng bao nhiêu? Chọn giá trị đúng.
\choice
{18.75}
{38.5}
{\True 37.5}
{75}
\loigiai{
Đáp án C. 37.5. Đối chiếu định nghĩa, hệ thức hoặc dữ kiện của bài toán cho kết luận trên.
}
\end{ex}

% ===== Câu 145 =====
% ID: L11C2B8-40-TL
% Mức: VD
\dangbt{Xác định bước sóng, chu kì, tần số và tốc độ truyền sóng}
\begin{ex}
% ID: L11C2B8-40-TL
% Mức: VD
Một học sinh giải bài sau nhưng bỏ qua điều kiện vật lí. Hãy sửa cách làm và trình bày lời giải đúng: Khoảng cách giữa đỉnh thứ nhất và đỉnh thứ mười một là $25\,\text{m}$. Sóng truyền với tốc độ $10\,\text{m}$ /s. Tính bước sóng và tần số.
\choiceTF

\loigiai{
Cách giải đúng: Giữa 11 đỉnh có 10 bước sóng nên $\lambda=25/10=2,5$ m; $f=v/\lambda=4$ Hz. Sai lầm cần tránh là dùng hệ thức khi chưa xác định điều kiện biên, môi trường hoặc đơn vị.
}
\end{ex}

% ===== Câu 146 =====
% ID: L11C2B8-40-TN
% Mức: VD
\dangbt{Phân tích trạng thái và chiều truyền sóng từ dữ kiện quan sát}
\begin{ex}
% ID: L11C2B8-40-TN
% Mức: VD
Trong chuyên đề “Phân tích trạng thái và chiều truyền sóng từ dữ kiện quan sát”, Kết luận nào phù hợp nhất khi giải bài toán sau: Chứng minh hai điểm cách nhau một số nguyên lần bước sóng dao động cùng pha.
\choice
{Không cần sử dụng định nghĩa hay dữ kiện đã cho.}
{Đại lượng cần tìm luôn bằng 0 trong mọi trường hợp.}
{Kết quả không phụ thuộc môi trường, tần số hoặc điều kiện biên.}
{\True Vì $\Delta\varphi=2\pi d/\lambda$. Với $d=k\lambda$ thì $\Delta\varphi=2k\pi$, nên hai điểm cùng pha.}
\loigiai{
Đáp án D. Vì $\Delta\varphi=2\pi d/\lambda$. Với $d=k\lambda$ thì $\Delta\varphi=2k\pi$, nên hai điểm cùng pha.. Đối chiếu định nghĩa, hệ thức hoặc dữ kiện của bài toán cho kết luận trên.
}
\end{ex}

% ===== Câu 147 =====
% ID: L11C2B8-41-TL
% Mức: VD
\dangbt{Xác định bước sóng, chu kì, tần số và tốc độ truyền sóng}
\begin{ex}
% ID: L11C2B8-41-TL
% Mức: VD
Mở rộng tình huống sau thành một ví dụ thực tế và trình bày phương pháp giải: Một sóng đi từ môi trường $A$ sang $B$, tần số $20\,\text{Hz}$ không đổi. Biết tốc độ ở $A$ là $40\,\text{m}$ /s, ở $B$ là $30\,\text{m}$ /s. So sánh bước sóng ở hai môi trường.
\choiceTF

\loigiai{
$\lambda_A=2$ m, $\lambda_B=1,5$ m; bước sóng ở $B$ nhỏ hơn ở $A0,5\,\text{m}$. Có thể kiểm chứng bằng cách thay kết quả vào dữ kiện ban đầu và đối chiếu với hiện tượng thực tế.
}
\end{ex}

% ===== Câu 148 =====
% ID: L11C2B8-41-TN
% Mức: VD
\dangbt{Phân tích trạng thái và chiều truyền sóng từ dữ kiện quan sát}
\begin{ex}
% ID: L11C2B8-41-TN
% Mức: VD
Trong chuyên đề “Phân tích trạng thái và chiều truyền sóng từ dữ kiện quan sát”, Kết luận nào phù hợp nhất khi giải bài toán sau: Một ảnh sóng có bước sóng $80\,\text{cm}$. Xác định ba vị trí gần nhất tính từ $O$ theo chiều truyền mà phần tử ngược pha với O.
\choice
{\True Điểm ngược pha thỏa $d=(k+1/2)\lambda$. Ba vị trí gần nhất là $40\,\text{cm}$, $120\,\text{cm}$ và $200\,\text{cm}$.}
{Không cần sử dụng định nghĩa hay dữ kiện đã cho.}
{Đại lượng cần tìm luôn bằng 0 trong mọi trường hợp.}
{Kết quả không phụ thuộc môi trường, tần số hoặc điều kiện biên.}
\loigiai{
Đáp án A. Điểm ngược pha thỏa $d=(k+1/2)\lambda$. Ba vị trí gần nhất là $40\,\text{cm}$, $120\,\text{cm}$ và $200\,\text{cm}$.. Đối chiếu định nghĩa, hệ thức hoặc dữ kiện của bài toán cho kết luận trên.
}
\end{ex}

% ===== Câu 149 =====
% ID: L11C2B8-42-TN
% Mức: VD
\begin{ex}
% ID: L11C2B8-42-TN
% Mức: VD
Trong chuyên đề “Phân tích trạng thái và chiều truyền sóng từ dữ kiện quan sát”, Kết luận nào phù hợp nhất khi giải bài toán sau: Giải thích vì sao đường cong hình sin tại một thời điểm không phải là quỹ đạo chuyển động của một phần tử môi trường.
\choice
{Kết quả không phụ thuộc môi trường, tần số hoặc điều kiện biên.}
{\True Đường cong biểu diễn li độ của nhiều phần tử ở các vị trí khác nhau tại cùng thời điểm}
{Không cần sử dụng định nghĩa hay dữ kiện đã cho.}
{Đại lượng cần tìm luôn bằng 0 trong mọi trường hợp.}
\loigiai{
Đáp án B. Đường cong biểu diễn li độ của nhiều phần tử ở các vị trí khác nhau tại cùng thời điểm. Đối chiếu định nghĩa, hệ thức hoặc dữ kiện của bài toán cho kết luận trên.
}
\end{ex}

% ===== Câu 150 =====
% ID: L11C2B8-43-TN
% Mức: VD
\begin{ex}
% ID: L11C2B8-43-TN
% Mức: VD
Trong chuyên đề “Phân tích trạng thái và chiều truyền sóng từ dữ kiện quan sát”, Phương pháp phù hợp nhất cho dạng “Phân tích trạng thái và chiều truyền sóng từ dữ kiện quan sát” là
\choice
{bỏ qua đơn vị và chiều truyền sóng}
{luôn coi mọi điểm dao động cùng pha}
{\True xác định dữ kiện, chọn đúng hệ thức hoặc định nghĩa, đổi đơn vị rồi kiểm tra kết quả}
{chỉ thay số mà không cần xét điều kiện áp dụng}
\loigiai{
Đáp án C. xác định dữ kiện, chọn đúng hệ thức hoặc định nghĩa, đổi đơn vị rồi kiểm tra kết quả. Đối chiếu định nghĩa, hệ thức hoặc dữ kiện của bài toán cho kết luận trên.
}
\end{ex}

% ===== Câu 151 =====
% ID: L11C2B8-44-TN
% Mức: NB
\dangbt{Vận dụng khái niệm sóng cơ trong các hiện tượng thực tế}
\begin{ex}
% ID: L11C2B8-44-TN
% Mức: NB
Trong chuyên đề “Vận dụng khái niệm sóng cơ trong các hiện tượng thực tế”, Khi sóng cơ truyền qua một môi trường, đại lượng được truyền đi từ nơi này đến nơi khác là
\choice
{các phần tử vật chất}
{khối lượng môi trường}
{điện tích của phần tử}
{\True năng lượng và trạng thái dao động}
\loigiai{
Đáp án D. năng lượng và trạng thái dao động. Đối chiếu định nghĩa, hệ thức hoặc dữ kiện của bài toán cho kết luận trên.
}
\end{ex}

% ===== Câu 152 =====
% ID: L11C2B8-45-TN
% Mức: NB
\begin{ex}
% ID: L11C2B8-45-TN
% Mức: NB
Trong chuyên đề “Vận dụng khái niệm sóng cơ trong các hiện tượng thực tế”, Sóng cơ không thể truyền trong
\choice
{\True chân không}
{thép}
{nước}
{không khí}
\loigiai{
Đáp án A. chân không. Đối chiếu định nghĩa, hệ thức hoặc dữ kiện của bài toán cho kết luận trên.
}
\end{ex}

% ===== Câu 153 =====
% ID: L11C2B8-46-TN
% Mức: TH
\begin{ex}
% ID: L11C2B8-46-TN
% Mức: TH
Trong chuyên đề “Vận dụng khái niệm sóng cơ trong các hiện tượng thực tế”, Một nút chai nổi trên mặt hồ khi có sóng truyền qua chủ yếu
\choice
{đứng yên tuyệt đối}
{\True dao động quanh vị trí cân bằng}
{trôi theo đúng tốc độ truyền sóng}
{chuyển thẳng đến bờ}
\loigiai{
Đáp án B. dao động quanh vị trí cân bằng. Đối chiếu định nghĩa, hệ thức hoặc dữ kiện của bài toán cho kết luận trên.
}
\end{ex}

% ===== Câu 154 =====
% ID: L11C2B8-47-TN
% Mức: TH
\begin{ex}
% ID: L11C2B8-47-TN
% Mức: TH
Trong chuyên đề “Vận dụng khái niệm sóng cơ trong các hiện tượng thực tế”, Một sóng truyền quãng đường $18\,\text{m}$ trong $6\,\text{s}$. Tốc độ truyền sóng tính theo m/s bằng bao nhiêu? Chọn giá trị đúng.
\choice
{01/05/2026}
{4}
{\True 3}
{6}
\loigiai{
Đáp án C. 3. Đối chiếu định nghĩa, hệ thức hoặc dữ kiện của bài toán cho kết luận trên.
}
\end{ex}

% ===== Câu 155 =====
% ID: L11C2B8-48-TN
% Mức: TH
\begin{ex}
% ID: L11C2B8-48-TN
% Mức: TH
Trong chuyên đề “Vận dụng khái niệm sóng cơ trong các hiện tượng thực tế”, Sóng truyền $24\,\text{m}$ trong $8\,\text{s}$. Tốc độ truyền sóng tính theo m/s bằng bao nhiêu? Chọn giá trị đúng.
\choice
{6}
{01/05/2026}
{4}
{\True 3}
\loigiai{
Đáp án D. 3. Đối chiếu định nghĩa, hệ thức hoặc dữ kiện của bài toán cho kết luận trên.
}
\end{ex}

% ===== Câu 156 =====
% ID: L11C2B8-49-TN
% Mức: TH
\begin{ex}
% ID: L11C2B8-49-TN
% Mức: TH
Trong chuyên đề “Vận dụng khái niệm sóng cơ trong các hiện tượng thực tế”, Một xung sóng đi được $15\,\text{m}$ với tốc độ $5\,\text{m}$ /s. Thời gian truyền tính theo s bằng bao nhiêu? Chọn giá trị đúng.
\choice
{\True 3}
{6}
{01/05/2026}
{4}
\loigiai{
Đáp án A. 3. Đối chiếu định nghĩa, hệ thức hoặc dữ kiện của bài toán cho kết luận trên.
}
\end{ex}

% ===== Câu 157 =====
% ID: L11C2B8-50-TN
% Mức: VD
\begin{ex}
% ID: L11C2B8-50-TN
% Mức: VD
Trong chuyên đề “Vận dụng khái niệm sóng cơ trong các hiện tượng thực tế”, Kết luận nào phù hợp nhất khi giải bài toán sau: Giải thích vì sao sóng cơ truyền năng lượng nhưng không mang các phần tử môi trường đi theo.
\choice
{Kết quả không phụ thuộc môi trường, tần số hoặc điều kiện biên.}
{\True Các phần tử môi trường chỉ dao động quanh vị trí cân bằng}
{Không cần sử dụng định nghĩa hay dữ kiện đã cho.}
{Đại lượng cần tìm luôn bằng 0 trong mọi trường hợp.}
\loigiai{
Đáp án B. Các phần tử môi trường chỉ dao động quanh vị trí cân bằng. Đối chiếu định nghĩa, hệ thức hoặc dữ kiện của bài toán cho kết luận trên.
}
\end{ex}

% ===== Câu 158 =====
% ID: L11C2B8-51-TN
% Mức: VD
\begin{ex}
% ID: L11C2B8-51-TN
% Mức: VD
Trong chuyên đề “Vận dụng khái niệm sóng cơ trong các hiện tượng thực tế”, Kết luận nào phù hợp nhất khi giải bài toán sau: Nêu hai điều kiện cần để hình thành sóng cơ từ một nguồn dao động.
\choice
{Đại lượng cần tìm luôn bằng 0 trong mọi trường hợp.}
{Kết quả không phụ thuộc môi trường, tần số hoặc điều kiện biên.}
{\True Cần có nguồn dao động và môi trường vật chất có lực liên kết giữa các phần tử để truyền dao động.}
{Không cần sử dụng định nghĩa hay dữ kiện đã cho.}
\loigiai{
Đáp án C. Cần có nguồn dao động và môi trường vật chất có lực liên kết giữa các phần tử để truyền dao động.. Đối chiếu định nghĩa, hệ thức hoặc dữ kiện của bài toán cho kết luận trên.
}
\end{ex}

% ===== Câu 159 =====
% ID: L11C2B8-52-TN
% Mức: VD
\begin{ex}
% ID: L11C2B8-52-TN
% Mức: VD
Trong chuyên đề “Vận dụng khái niệm sóng cơ trong các hiện tượng thực tế”, Kết luận nào phù hợp nhất khi giải bài toán sau: Giải thích vì sao không thể nghe trực tiếp tiếng chuông đặt trong một bình đã hút hết không khí.
\choice
{Không cần sử dụng định nghĩa hay dữ kiện đã cho.}
{Đại lượng cần tìm luôn bằng 0 trong mọi trường hợp.}
{Kết quả không phụ thuộc môi trường, tần số hoặc điều kiện biên.}
{\True Âm là sóng cơ nên cần môi trường vật chất để truyền}
\loigiai{
Đáp án D. Âm là sóng cơ nên cần môi trường vật chất để truyền. Đối chiếu định nghĩa, hệ thức hoặc dữ kiện của bài toán cho kết luận trên.
}
\end{ex}

% ===== Câu 160 =====
% ID: L11C2B8-53-TN
% Mức: VD
\begin{ex}
% ID: L11C2B8-53-TN
% Mức: VD
Trong chuyên đề “Vận dụng khái niệm sóng cơ trong các hiện tượng thực tế”, Phương pháp phù hợp nhất cho dạng “Vận dụng khái niệm sóng cơ trong các hiện tượng thực tế” là
\choice
{\True xác định dữ kiện, chọn đúng hệ thức hoặc định nghĩa, đổi đơn vị rồi kiểm tra kết quả}
{chỉ thay số mà không cần xét điều kiện áp dụng}
{bỏ qua đơn vị và chiều truyền sóng}
{luôn coi mọi điểm dao động cùng pha}
\loigiai{
Đáp án A. xác định dữ kiện, chọn đúng hệ thức hoặc định nghĩa, đổi đơn vị rồi kiểm tra kết quả. Đối chiếu định nghĩa, hệ thức hoặc dữ kiện của bài toán cho kết luận trên.
}
\end{ex}

% ===== Câu 161 =====
% ID: L11C2B8-54-TN
% Mức: NB
\dangbt{Xác định bước sóng, chu kì, tần số và tốc độ truyền sóng}
\begin{ex}
% ID: L11C2B8-54-TN
% Mức: NB
Một sóng có tần số $5\,\text{Hz}$ và tốc độ $20\,\text{m}$ /s. Bước sóng bằng
\choice
{$15\,\text{m}$}
{\True $4\,\text{m}$}
{$100\,\text{m}$}
{$0,25\,\text{m}$}
\loigiai{
Đáp án B. $4\,\text{m}$. Đối chiếu định nghĩa, hệ thức hoặc dữ kiện của bài toán cho kết luận trên.
}
\end{ex}

% ===== Câu 162 =====
% ID: L11C2B8-55-TN
% Mức: NB
\begin{ex}
% ID: L11C2B8-55-TN
% Mức: NB
Khoảng cách giữa 7 đỉnh sóng liên tiếp là $12\,\text{m}$. Bước sóng bằng
\choice
{$6\,\text{m}$}
{$84\,\text{m}$}
{\True $2\,\text{m}$}
{12/ $7\,\text{m}$}
\loigiai{
Đáp án C. $2\,\text{m}$. Đối chiếu định nghĩa, hệ thức hoặc dữ kiện của bài toán cho kết luận trên.
}
\end{ex}

% ===== Câu 163 =====
% ID: L11C2B8-56-TN
% Mức: TH
\begin{ex}
% ID: L11C2B8-56-TN
% Mức: TH
Sóng truyền với bước sóng $0,8\,\text{m}$ và chu kì $0,2\,\text{s}$. Tốc độ truyền sóng bằng
\choice
{$0,16\,\text{m}$ /s}
{$1\,\text{m}$ /s}
{$16\,\text{m}$ /s}
{\True $4\,\text{m}$ /s}
\loigiai{
Đáp án D. $4\,\text{m}$ /s. Đối chiếu định nghĩa, hệ thức hoặc dữ kiện của bài toán cho kết luận trên.
}
\end{ex}

% ===== Câu 164 =====
% ID: L11C2B8-57-TN
% Mức: TH
\begin{ex}
% ID: L11C2B8-57-TN
% Mức: TH
Sóng có tốc độ $12\,\text{m}$ /s và tần số $4\,\text{Hz}$. Bước sóng tính theo m bằng bao nhiêu? Chọn giá trị đúng.
\choice
{\True 3}
{5.1}
{01/05/2026}
{4}
\loigiai{
Đáp án A. 3. Đối chiếu định nghĩa, hệ thức hoặc dữ kiện của bài toán cho kết luận trên.
}
\end{ex}

% ===== Câu 165 =====
% ID: L11C2B8-58-TN
% Mức: TH
\begin{ex}
% ID: L11C2B8-58-TN
% Mức: TH
Khoảng cách giữa 9 đỉnh sóng liên tiếp là $16\,\text{m}$. Bước sóng tính theo m bằng bao nhiêu? Chọn giá trị đúng.
\choice
{3}
{\True 2}
{4}
{1}
\loigiai{
Đáp án B. 2. Đối chiếu định nghĩa, hệ thức hoặc dữ kiện của bài toán cho kết luận trên.
}
\end{ex}

% ===== Câu 166 =====
% ID: L11C2B8-59-TN
% Mức: TH
\begin{ex}
% ID: L11C2B8-59-TN
% Mức: TH
Sóng có bước sóng $1,5\,\text{m}$ và chu kì $0,25\,\text{s}$. Tốc độ truyền tính theo m/s bằng bao nhiêu? Chọn giá trị đúng.
\choice
{3}
{7}
{\True 6}
{12}
\loigiai{
Đáp án C. 6. Đối chiếu định nghĩa, hệ thức hoặc dữ kiện của bài toán cho kết luận trên.
}
\end{ex}

% ===== Câu 167 =====
% ID: L11C2B8-60-TN
% Mức: VD
\begin{ex}
% ID: L11C2B8-60-TN
% Mức: VD
Kết luận nào phù hợp nhất khi giải bài toán sau: Một sóng có tần số $8\,\text{Hz}$ truyền với tốc độ $24\,\text{m}$ /s. Tính chu kì, bước sóng và quãng đường sóng truyền trong $1,5\,\text{s}$.
\choice
{Không cần sử dụng định nghĩa hay dữ kiện đã cho.}
{Đại lượng cần tìm luôn bằng 0 trong mọi trường hợp.}
{Kết quả không phụ thuộc môi trường, tần số hoặc điều kiện biên.}
{\True $T=1/f=0,125$ s}
\loigiai{
Đáp án D. $T=1/f=0,125$ s. Đối chiếu định nghĩa, hệ thức hoặc dữ kiện của bài toán cho kết luận trên.
}
\end{ex}

% ===== Câu 168 =====
% ID: L11C2B8-61-TN
% Mức: VD
\begin{ex}
% ID: L11C2B8-61-TN
% Mức: VD
Kết luận nào phù hợp nhất khi giải bài toán sau: Khoảng cách giữa đỉnh thứ nhất và đỉnh thứ mười một là $25\,\text{m}$. Sóng truyền với tốc độ $10\,\text{m}$ /s. Tính bước sóng và tần số.
\choice
{\True Giữa 11 đỉnh có 10 bước sóng nên $\lambda=25/10=2,5$ m}
{Không cần sử dụng định nghĩa hay dữ kiện đã cho.}
{Đại lượng cần tìm luôn bằng 0 trong mọi trường hợp.}
{Kết quả không phụ thuộc môi trường, tần số hoặc điều kiện biên.}
\loigiai{
Đáp án A. Giữa 11 đỉnh có 10 bước sóng nên $\lambda=25/10=2,5$ m. Đối chiếu định nghĩa, hệ thức hoặc dữ kiện của bài toán cho kết luận trên.
}
\end{ex}

% ===== Câu 169 =====
% ID: L11C2B8-62-TN
% Mức: VD
\begin{ex}
% ID: L11C2B8-62-TN
% Mức: VD
Kết luận nào phù hợp nhất khi giải bài toán sau: Một sóng đi từ môi trường $A$ sang $B$, tần số $20\,\text{Hz}$ không đổi. Biết tốc độ ở $A$ là $40\,\text{m}$ /s, ở $B$ là $30\,\text{m}$ /s. So sánh bước sóng ở hai môi trường.
\choice
{Kết quả không phụ thuộc môi trường, tần số hoặc điều kiện biên.}
{\True $\lambda_A=2$ m, $\lambda_B=1,5$ m}
{Không cần sử dụng định nghĩa hay dữ kiện đã cho.}
{Đại lượng cần tìm luôn bằng 0 trong mọi trường hợp.}
\loigiai{
Đáp án B. $\lambda_A=2$ m, $\lambda_B=1,5$ m. Đối chiếu định nghĩa, hệ thức hoặc dữ kiện của bài toán cho kết luận trên.
}
\end{ex}

% ===== Câu 170 =====
% ID: L11C2B8-63-TN
% Mức: VD
\begin{ex}
% ID: L11C2B8-63-TN
% Mức: VD
Phương pháp phù hợp nhất cho dạng “Xác định bước sóng, chu kì, tần số và tốc độ truyền sóng” là
\choice
{bỏ qua đơn vị và chiều truyền sóng}
{luôn coi mọi điểm dao động cùng pha}
{\True xác định dữ kiện, chọn đúng hệ thức hoặc định nghĩa, đổi đơn vị rồi kiểm tra kết quả}
{chỉ thay số mà không cần xét điều kiện áp dụng}
\loigiai{
Đáp án C. xác định dữ kiện, chọn đúng hệ thức hoặc định nghĩa, đổi đơn vị rồi kiểm tra kết quả. Đối chiếu định nghĩa, hệ thức hoặc dữ kiện của bài toán cho kết luận trên.
}
\end{ex}
